% 第 1 章 统一书稿前言包装
\newcommand{\lectureno}{第一章}
\newcommand{\lectureshort}{从可行集到模型价值}
% ============================================================
%  微观经济学理论课程讲义 —— 统一课程前言（四讲共用）
%  版式参考 Nolan Miller, Notes on Microeconomic Theory (SWP/notes3 风格)
%  编译：LuaLaTeX（两遍）
%  各讲 preamble.tex 只需设定 \lectureno / \lectureshort 后 \input 本文件。
% ============================================================
\documentclass[11pt,a4paper,fontset=none,oneside,openany]{ctexbook}

% ---------- 版式 ----------
\usepackage{geometry}
\geometry{
  a4paper,
  top=2.8cm, bottom=2.8cm,
  left=2.6cm, right=2.6cm,
  headsep=0.85cm, footskip=1.25cm
}
\setlength{\headheight}{14pt}

% ---------- 字体（全部矢量嵌入，保证 PDF 高清晰度）----------
% 注意：本机 macOS 已移除 Kaiti SC，CJK 斜体统一回退到 Songti SC；
%        Songti SC / Heiti SC 经 OSFONTDIR 指向 Supplemental 后由 luaotfload 解析。
\usepackage{fontspec}
\setmainfont{TeX Gyre Termes}[Ligatures=TeX]
\setsansfont{TeX Gyre Heros}[Ligatures=TeX,Scale=MatchLowercase]
\setmonofont{TeX Gyre Cursor}[Scale=MatchLowercase]
\setCJKmainfont{Songti SC}[BoldFont={Songti SC Bold},ItalicFont={Songti SC}]
\setCJKsansfont{Heiti SC}[BoldFont={Heiti SC}]

% ---------- 数学 ----------
\usepackage{amsmath}
\usepackage{mathtools}
\usepackage{amsthm}
\usepackage{unicode-math}
\setmathfont{TeX Gyre Termes Math}
\allowdisplaybreaks[2]
\setlength{\abovedisplayskip}{9pt plus 3pt minus 3pt}
\setlength{\belowdisplayskip}{9pt plus 3pt minus 3pt}

% ---------- 定理类环境（书稿格式：按章编号，如“定理 1.1”）----------
\theoremstyle{plain}
\newtheorem{theorem}{定理}[chapter]
\newtheorem{proposition}[theorem]{命题}
\newtheorem{lemma}[theorem]{引理}
\newtheorem{corollary}[theorem]{推论}
\theoremstyle{definition}
\newtheorem{definition}[theorem]{定义}
\newtheorem{condition}[theorem]{条件}
\newtheorem{example}[theorem]{例}
\theoremstyle{remark}
\newtheorem{remark}[theorem]{注}
% 公式、图、表按章编号（书稿格式：式 (1.1)、图 1.1、表 1.1）
\numberwithin{equation}{chapter}
\numberwithin{figure}{chapter}
\numberwithin{table}{chapter}

% 定理类环境标题后的标点：amsthm 默认加英文 "."，此处统一改为中文句号。
% 原理：amsthm 的 \@thm 每次开环境都执行 \thm@headpunct{.}，而样式宏在其之后展开，
% 故把 \thm@headpunct{。} 追加到三个样式宏末尾即可覆盖（用 \AtBeginDocument 以免被 ctex 覆盖）。
\makeatletter
\AtBeginDocument{%
  \let\WBth@plain\th@plain
  \let\WBth@definition\th@definition
  \let\WBth@remark\th@remark
  \def\th@plain{\WBth@plain\thm@headpunct{。}}%
  \def\th@definition{\WBth@definition\thm@headpunct{。}}%
  \def\th@remark{\WBth@remark\thm@headpunct{。}}%
}
\makeatother

% 无编号小标题（仿 notes3 的 “A Note on ...”）
\newcommand{\notehead}[1]{\par\medskip\noindent{\bfseries #1}\par\smallskip}

% ---------- 表格 ----------
\usepackage{booktabs}
\usepackage{array}
\renewcommand{\arraystretch}{1.25}
\setlength{\tabcolsep}{8pt}

% ---------- 图形 ----------
\usepackage{graphicx}
\usepackage{tikz}
\usepackage{pgfplots}
\pgfplotsset{compat=1.18}
\usetikzlibrary{arrows.meta,positioning,shapes.geometric,calc,patterns,decorations.pathreplacing,decorations.markings}
\usepackage{float}
\usepackage{caption}
\captionsetup{font=small,labelfont=bf,labelsep=quad,skip=4pt}

% 统一图形配色（黑白模式：黑/灰梯度，供所有 TikZ 图复用）
% —— 全局黑白化：原彩色调色板统一改为黑/灰，保证打印与投影均可读 ——
\definecolor{AxisBlue}{RGB}{0,0,0}        % 轴线 / 主曲线：黑
\definecolor{SoftBlue}{RGB}{130,130,130}  % 节点填充：中灰
\definecolor{LightRed}{RGB}{224,224,224}  % 浅填充：浅灰
\definecolor{DeepBlue}{RGB}{0,0,0}        % 描边：黑
\definecolor{MidBlue}{RGB}{70,70,70}      % 次曲线 / 网格：深灰
\tikzset{
  >=Stealth,
  every picture/.append style={font=\small},
  axisline/.style={thick,color=black,-Stealth},
  axislabel/.style={color=black},
  fillsoft/.style={fill=LightRed,fill opacity=0.6,draw=black,thick},
  nodeblue/.style={circle,fill=black!85,draw=black,thick,inner sep=2pt},
  nodewhite/.style={circle,fill=white,draw=black,thick,inner sep=2pt},
  arrblue/.style={->, thick, color=black},
  boxsoft/.style={draw=black, rounded corners=2pt, align=center, inner sep=5pt, fill=black!6},
  % 多曲线 / 多集合区分：以线型 + 标记代替颜色
  curveA/.style={thick,black,solid},
  curveB/.style={thick,black,dashed},
  curveC/.style={thick,black,dotted},
  markA/.style={mark=*,mark size=2.2pt,color=black},
  markB/.style={mark=square*,mark size=2.0pt,color=black},
  markC/.style={mark=triangle*,mark size=2.4pt,color=black},
  % 图内区域 / 点标签（白底避免与图形重叠）
  figlabel/.style={font=\scriptsize,fill=white,inner sep=1pt,outer sep=1pt,text=black},
}
% 图形说明（黑白模式下对图形区域 / 曲线 / 关键点的解释，置于图形下方）
\newcommand{\figexplan}[1]{%
  \par\smallskip\noindent{\small\itshape 图形说明：#1}\par\smallskip}

% ---------- 页眉页脚（书稿格式：左=书名，右=章次与章名）----------
\usepackage{fancyhdr}
% 章标题：居中、大号（书稿章首样式）
\CTEXsetup[format={\centering\LARGE\bfseries}]{chapter}
\CTEXsetup[beforeskip=16pt,afterskip=14pt]{chapter}
\CTEXsetup[format={\large\bfseries}]{section}
\CTEXsetup[format={\normalsize\bfseries}]{subsection}
\pagestyle{fancy}
\fancyhf{}
\providecommand{\booktitle}{微观经济学理论：可行集、偏好与效用}
\providecommand{\lectureno}{第一章}
\providecommand{\lectureshort}{课程讲义}
\fancyhead[L]{\small\itshape \booktitle}
\fancyhead[R]{\small \lectureno\quad \lectureshort}
\fancyfoot[C]{\small \thepage}
\renewcommand{\headrulewidth}{0.4pt}
\fancypagestyle{plain}{\fancyhf{}\fancyfoot[C]{\small \thepage}\renewcommand{\headrulewidth}{0pt}}

% ---------- 正文细节 ----------
\usepackage{setspace}
\setstretch{1.12}
\usepackage{microtype}
\tolerance=900
\emergencystretch=2em
\raggedbottom

% ---------- 超链接与 PDF 书签 ----------
\usepackage{hyperref}
\hypersetup{
  colorlinks=true,
  linkcolor=black,
  urlcolor=black,
  citecolor=black,
  pdfencoding=auto,
  bookmarks=true,
  pdftitle={\booktitle\quad \lectureno：\lectureshort},
  pdfsubject={微观经济学理论书稿}
}
% 章节标题中常见数学符号的安全书签映射（避免 "Improper alphabetic constant"）
\pdfstringdefDisableCommands{%
  \def\leftrightarrow{<->}%
  \def\star{*}%
  \def\S{S}%
  \def\succeq{\text{≽}}%
  \def\succ{\text{≻}}%
  \def\sim{\text{∼}}%
}

% ---------- 常用记号 ----------
\newcommand{\R}{\mathbb{R}}
\newcommand{\N}{\mathbb{N}}
\newcommand{\Z}{\mathbb{Z}}
\newcommand{\Q}{\mathbb{Q}}
\newcommand{\succeqq}{\succeq}
\newcommand{\argmax}{\operatorname*{arg\,max}}
\newcommand{\argmin}{\operatorname*{arg\,min}}
\newcommand{\st}{\text{s.t.}\ }

\setcounter{chapter}{0}


\begin{document}

\chapter{从可行集到模型价值：选择理论的形式化基础与方法论边界}
\noindent{\normalsize\itshape Individual decision-making under certainty\par}
\vspace{0.5cm}

\section*{本章定位与学习目标}

本章从确定性下的个体决策出发，建立一条严格的形式化链条：
\[
\theta \;\longmapsto\; B(\theta) \;\longmapsto\; C(\theta)
\;\longmapsto\; \text{行为预测} \;\longmapsto\; \text{Positive / Normative analysis}.
\]

完成本章后，读者应能够：
\begin{enumerate}
  \item 区分方案空间、可行集、函数与对应，并准确使用集合语言；
  \item 解释为何选择对应的输出是\emph{集合}而非点；
  \item 准确使用
        \[
          C(\theta)=\argmax\{\,u(x):x\in B(\theta)\,\};
        \]
  \item 说明纯选择理论为何可以不预设目标函数；
  \item 区分正向命题（positive statement）与规范判断（normative statement）；
  \item 以有用性、准确性、广泛性、简洁性、识别性与可证伪性评价一个模型。
\end{enumerate}

\notehead{本章内容与主线}

\begin{center}
\begin{tabular}{clp{6.5cm}}
\toprule
次序 & 单元 & 核心任务 \\
\midrule
1 & 导入 & 从现实选择到形式化对象；提出“模型究竟在映射什么？” \\
2 & 集合与可行集 & 集合语言、实数空间、方案空间、预算集 \\
3 & 函数与对应 & 单值函数与集值映射；为什么输出可以是集合 \\
4 & 参数与选择 & 预算集、$\argmax$、最大值函数、比较静态 \\
5 & 纯选择理论 & 不预设目标函数；理性化与显示偏好 \\
6 & Positive / Normative & 事实、因果、价值与福利的三层分离 \\
7 & 模型价值 & Useful、Accurate、Broad、Compact、Identifiable、Falsifiable \\
8 & 本章要点 & 一条逻辑链、三个反例、一道综合题 \\
\bottomrule
\end{tabular}
\end{center}

\notehead{本章核心公式链}
\[
X\supseteq B(\theta)\;;\qquad
u:X\to\R\;;\qquad
C(\theta)\subseteq B(\theta)\;;\qquad
C(\theta)=\argmax_{x\in B(\theta)} u(x).
\]
随后写：$v(\theta)=\max u$ 是\emph{数}，$C(\theta)$ 是\emph{集合}；最后写：Positive 是事实与机制，Normative 是价值与评价。

\clearpage
\tableofcontents
\clearpage

% ==================== 第 1 节：导入（0–10 分钟）====================
\section{导入：从现实选择到形式化对象}

\notehead{1.1 一个真实的决策问题}

想象一位博士生每月有可支配收入 $m=3000$ 元，需要在两类支出之间分配：书籍（每单位价格 $p_1=100$ 元）与餐饮（每单位价格 $p_2=50$ 元）。她能买“任意多”吗？不能。她买得起哪些组合？她最终会选哪个组合？我们能不能\emph{预测}她的选择、又能不能\emph{评价}她的选择“好不好”？

把这三个问题逐一形式化，就是本章的全部内容：
\begin{enumerate}
  \item \textbf{买得起哪些} $\;\Longrightarrow\;$ 一个\emph{集合}问题（可行集）；
  \item \textbf{会选哪个} $\;\Longrightarrow\;$ 一个\emph{映射}问题（选择对应）；
  \item \textbf{好不好} $\;\Longrightarrow\;$ 一个\emph{方法论}问题（实证与规范的分界）。
\end{enumerate}

\begin{figure}[H]
\centering
\begin{tikzpicture}[font=\small,scale=1.12,transform shape]
  % 左面板：现实语言
  \draw[thick,fill=gray!15] (0,0) rectangle (3.4,3.0);
  \node[font=\bfseries] at (1.7,2.62) {现实语言};
  \node[align=center,text width=3.1cm,font=\scriptsize] at (1.7,1.82) {月收入 3000 元；书籍单价 100 元\\餐饮单价 50 元；“这本书买不买？”};
  \node[align=center,font=\scriptsize] at (1.7,0.62) {无法直接计算：\\既没有集合，也没有映射};
  \node[font=\scriptsize] at (1.7,-0.48) {(a) 现实决策情境};
  % 形式化箭头
  \draw[-{Stealth[length=2.6mm]},thick] (3.6,1.5) -- (4.4,1.5);
  \node[font=\scriptsize,above=3pt] at (4.0,1.5) {形式化};
  % 右面板：形式语言
  \begin{scope}[xshift=8.0cm]
    \draw[thick,fill=gray!15] (0,0) rectangle (3.4,3.0);
    \node[font=\bfseries] at (1.7,2.62) {形式语言};
    \node[align=center,text width=3.1cm,font=\scriptsize] at (1.7,1.82) {$B(p,m)$：预算集（集合）\\$C(\theta)$：选择对应（映射）\\实证与规范之分（方法论）};
    \node[align=center,font=\scriptsize] at (1.7,0.62) {可计算、可检验、可反驳};
    \node[font=\scriptsize] at (1.7,-0.48) {(b) 形式化对象};
  \end{scope}
\end{tikzpicture}
\caption{从现实决策到形式化对象：同一问题的两种语言}
\end{figure}
\figexplan{左面板是被提问的真实情境，其中没有集合、也没有映射，因此无法直接计算；经过中间的形式化箭头，右面板把同一问题拆成三个可操作的对象：预算集 $B(p,m)$ 是集合、选择对应 $C(\theta)$ 是映射、实证与规范之别是方法论区分。三者正是后续各节依次处理的对象。}

\begin{example}[预算约束的形式化]
给定收入 $m=3000$，价格 $p_1=100$、$p_2=50$，可行集为
\[
B(p,m)=\{(x_1,x_2)\in\R_+^2:100x_1+50x_2\le 3000\}.
\]
预算线为 $100x_1+50x_2=3000$，即 $x_2=60-2x_1$；在 $x_1$ 轴截距为 30，在 $x_2$ 轴截距为 60。若该生选择 10 本书，则最多可支配 $(3000-1000)/50=40$ 单位餐饮。选择哪一个 $(x_1,x_2)$，是第 4--5 节的问题；判断该选择好不好，是第 6 节的问题。
\end{example}

\notehead{1.2 模型究竟在映射什么？}

理论经济学的第一步不是写方程，而是回答：“模型里的每个符号，对应现实世界里的什么？”本章的形式化链条是
\[
\theta \;\longmapsto\; B(\theta) \;\longmapsto\; C(\theta)
\;\longmapsto\; \text{行为预测} \;\longmapsto\; \text{Positive / Normative analysis},
\]
其中每一环都是一次\emph{抽象}：
\begin{itemize}
  \item $\theta$：模型外的给定条件（价格、收入、税率、禀赋……），我们称之为\emph{参数}；
  \item $B(\theta)$：在参数 $\theta$ 下，决策者可行的全部方案所成的\emph{集合}；
  \item $C(\theta)$：在可行集中，决策者\emph{实际会选}的方案所成的集合——注意它仍是集合；
  \item 行为预测：当 $\theta$ 变化时，$C(\theta)$ 如何变化（比较静态）；
  \item Positive / Normative：对预测结果分别作\emph{事实陈述}与\emph{价值评价}。
\end{itemize}

\medskip
\noindent\textbf{课堂提问：}在这个链条里，哪一环是“客观的”、哪一环可能掺入我们的价值判断？先记下你的直觉，第 6 节再回来检验。

% ==================== 第 2 节：集合与可行集（10–30 分钟）====================
\section{集合论基础与方案空间}

整个选择理论只用两种数学对象：\emph{集合}与\emph{映射}。本节先把集合语言讲透，再用它定义方案空间与可行集。

\subsection{集合的基本语言}

\begin{definition}[集合与元素]
\emph{集合}（set）是若干对象的总体，这些对象称为它的\emph{元素}。若 $a$ 是集合 $A$ 的元素，记 $a\in A$；否则记 $a\notin A$。
\end{definition}

集合有两种记法。\textbf{列举法}：$A=\{a,b,c\}$。\textbf{描述法}（刻画法）：
\[
A=\{\,x : \text{“}x\text{ 满足的性质”}\,\},
\qquad\text{例如}\qquad
E=\{\,n\in\Z : n \text{ 是偶数}\,\}.
\]

\begin{definition}[子集与相等]
若 $A$ 的每个元素都属于 $B$，称 $A$ 是 $B$ 的\emph{子集}，记 $A\subseteq B$。若 $A\subseteq B$ 且 $B\subseteq A$，则 $A=B$。若 $A\subseteq B$ 且 $A\neq B$，记 $A\subsetneq B$（真子集）。不含任何元素的集合称为\emph{空集}，记 $\varnothing$；对任意集合 $A$ 都有 $\varnothing\subseteq A$。
\end{definition}

证明两个集合相等的标准方法就是\textbf{双向包含}：先证 $A\subseteq B$，再证 $B\subseteq A$。这是本课乃至全部微观理论中最常用的证明套路。

\begin{definition}[并、交、差、补]
给定全集 $U$ 与 $A,B\subseteq U$：
\[
\begin{aligned}
A\cup B &= \{\,x : x\in A \ \text{或}\ x\in B\,\},\\
A\cap B &= \{\,x : x\in A \ \text{且}\ x\in B\,\},\\
A\setminus B &= \{\,x : x\in A \ \text{且}\ x\notin B\,\},\\
A^{c} &= U\setminus A .
\end{aligned}
\]
\end{definition}

\begin{proposition}[德摩根律]
$(A\cup B)^{c}=A^{c}\cap B^{c}$，\quad $(A\cap B)^{c}=A^{c}\cup B^{c}$。
\end{proposition}

\begin{proof}
只证第一式，第二式同理。先证 $(A\cup B)^{c}\subseteq A^{c}\cap B^{c}$：任取 $x\in(A\cup B)^{c}$，则 $x\notin A\cup B$，故 $x\notin A$ 且 $x\notin B$，即 $x\in A^{c}$ 且 $x\in B^{c}$，所以 $x\in A^{c}\cap B^{c}$。再证 $A^{c}\cap B^{c}\subseteq (A\cup B)^{c}$：任取 $x\in A^{c}\cap B^{c}$，则 $x\notin A$ 且 $x\notin B$，于是 $x\notin A\cup B$，即 $x\in(A\cup B)^{c}$。双向包含成立，故两集合相等。
\end{proof}

\begin{definition}[笛卡尔积]
集合 $A$ 与 $B$ 的\emph{笛卡尔积}是有序对的集合
\[
A\times B=\{\,(a,b):a\in A,\ b\in B\,\}.
\]
一般地，$A_1\times\cdots\times A_L=\{(a_1,\dots,a_L):a_l\in A_l\}$，其元素称为 $L$ 维\emph{有序数组}（向量）。
\end{definition}

\begin{example}[商品空间就是笛卡尔积]
设第 $l$ 种商品的消费量为 $x_l\in\R_+$（$\R_+$ 表示非负实数集），则一个消费束 $x=(x_1,\dots,x_L)$ 正是
\[
\R_+^{L}=\underbrace{\R_+\times\cdots\times\R_+}_{L\text{ 个}}
\]
中的一个点。当 $L=2$ 时，$\R_+^2$ 就是平面的第一象限（含坐标轴）。
\end{example}

\subsection{实数空间与几类常用集合}

微观经济学几乎全部工作在欧氏空间 $\R^{L}$ 里。以下记号务必牢记：
\[
\R_+=\{x\in\R:x\ge 0\},\qquad
\R_{++}=\{x\in\R:x>0\},\qquad
\R_+^{L}=\{x\in\R^{L}:x_l\ge 0,\ l=1,\dots,L\}.
\]

\begin{definition}[开球、开集与闭集]
在 $\R^{L}$ 中取范数 $\lVert x\rVert=\sqrt{\sum_{l}x_l^{2}}$。以 $x$ 为中心、$\varepsilon>0$ 为半径的\emph{开球}为
\[
N_{\varepsilon}(x)=\{\,y\in\R^{L}:\lVert y-x\rVert<\varepsilon\,\}.
\]
集合 $A$ 称为\emph{开集}，若其中每个点都有某个开球完全含于 $A$；$A$ 称为\emph{闭集}，若其补集 $\R^{L}\setminus A$ 是开集（等价地，$A$ 包含其全部极限点）。
\end{definition}

\begin{definition}[有界集与紧集]
$A\subseteq\R^{L}$ 称为\emph{有界}，若存在 $M>0$ 使一切 $x\in A$ 满足 $\lVert x\rVert\le M$。在 $\R^{L}$ 中，既是闭集又有界的集合称为\emph{紧集}（compact set）。
\end{definition}

\begin{remark}
“闭且有界 $\Rightarrow$ 紧”是有限维欧氏空间特有的 Heine–Borel 定理，在一般度量空间中不成立。本课只使用 $\R^{L}$ 上的版本。
\end{remark}

\begin{definition}[凸集]
$A\subseteq\R^{L}$ 称为\emph{凸集}，若对任意 $x,y\in A$ 与任意 $\lambda\in[0,1]$，都有
\[
\lambda x+(1-\lambda)y\in A.
\]
即连接集合中任意两点的线段仍全部落在集合内。
\end{definition}

\begin{figure}[H]
\centering
\begin{tikzpicture}[scale=1.0]
  % 凸集
  \draw[thick,fill=gray!15] (0,0) ellipse (1.5 and 1.0);
  \fill (0.5,0.4) circle (1.5pt) node[above] {\small $x$};
  \fill (-0.6,-0.5) circle (1.5pt) node[below] {\small $y$};
  \draw[thick] (0.5,0.4) -- (-0.6,-0.5);
  \node at (0,-1.6) {\small 凸集：线段仍在集内};
  % 非凸集
  \begin{scope}[xshift=6.2cm]
    \draw[thick,fill=gray!15]
      (-1.5,-1.0) .. controls (-1.8,0.6) and (-0.9,1.2) .. (0,0.9)
      .. controls (0.9,1.2) and (1.8,0.6) .. (1.5,-1.0)
      .. controls (0.6,-0.5) and (-0.6,-0.5) .. (-1.5,-1.0) -- cycle;
    \fill (-0.9,0.75) circle (1.5pt) node[above] {\small $x$};
    \fill (0.9,0.75) circle (1.5pt) node[above] {\small $y$};
    \draw[thick,dashed] (-0.9,0.75) -- (0.9,0.75);
    \node at (0,-1.6) {\small 非凸集：连线跑出集外};
  \end{scope}
\end{tikzpicture}
\caption{凸集与非凸集}
\end{figure}
\figexplan{左图为凸集（椭圆）：任取其中两点 $x,y$，连接它们的线段仍完全落在集合内部。右图为非凸集：同样取两点 $x,y$，连线（虚线）跑出集合之外。凸性保证可行集与偏好配合时最优解“长得规整”，是微观理论的核心工作假设。}

凸性是微观理论的“工作假设”：偏好凸保证无差异曲线内弯、技术凸保证边际报酬递减、可行集凸配合极值定理保证最优解“长得规整”。

\subsection{紧性：为什么“最优解存在”不是理所当然的}

前面的讨论已经说明：偏好可由连续效用表示、预算集是闭且有界的。但“闭且有界”并不是技术性的细节——它正是极值定理得以成立的前提。本小节把这件事讲透。

\begin{definition}[紧集（实数空间的刻画）]
设 $K\subseteq\R^{L}$。若 $K$ 同时是\emph{闭}的与\emph{有界}的，则称 $K$ 是\emph{紧}（compact）的。
\end{definition}

\begin{theorem}[Heine--Borel 定理]\label{thm:heineborel}
对 $K\subseteq\R^{L}$，以下三条等价：(i) $K$ 是紧的；(ii) $K$ 是\emph{闭且有界}的；
(iii) $K$ 中任何无穷序列都有收敛于 $K$ 的子列。
\end{theorem}

定理~\ref{thm:heineborel} 的第三条等价刻画最为有用：它把“紧”解释为“不能逃向无穷远”。反过来说，若集合不紧，就存在一个\emph{逃离序列}——序列中每一步都留在集合内，却越走越远，最终在极限处“跑到集合之外”。经济学中所有“最优解不存在”的反例，几乎都是这种逃离。

\begin{theorem}[分离定理（严格形式）]\label{thm:separate}
设 $A,B\subseteq\R^{L}$ 是两个\emph{不相交}的\emph{凸}集，且 $A$ 闭、$B$ 紧。则存在向量 $q\neq 0$ 与常数 $\alpha<\beta$ 使
\[
q\cdot a\ \le\ \alpha\ <\ \beta\ \le\ q\cdot b\qquad
\text{对一切 }a\in A,\ b\in B.
\]
\end{theorem}

定理~\ref{thm:separate} 是后续多处论证的引擎：凸偏好与凸预算集的分离、无差异曲线对最优束的分离、单调性与序贯选择的分离，背后都是它。其证明思路是：考虑函数 $\varphi(a,b)=\|a-b\|^{2}$ 在 $A\times B$ 上的取值；由 $B$ 紧且 $B$ 闭，最小值 $d^{2}>0$ 在某对 $(a^{*},b^{*})$ 处取到；取 $q=b^{*}-a^{*}$ 即可分离。

\begin{proposition}[预算集是紧集]\label{prop:budgetcompact}
设 $p\in\R_{++}^{L}$、$m\ge 0$，则 $B(p,m)=\{x\in\R^{L}_{+}:p\cdot x\le m\}$ 是紧集；$m>0$ 时它还满足“$x\in B(p,m)$ 蕴含 $\lambda x\in B(p,m)$ 对一切 $\lambda\in[0,1]$”，即\emph{向下闭合}。
\end{proposition}

\begin{proof}
$B(p,m)$ 显然有界：任取 $i$ 轴截距 $m/p_i$，由 $p\cdot x\le m$ 得 $0\le x_i\le m/p_i$。又 $B(p,m)$ 是有限个半空间 $\{x:p\cdot x\le m\}$ 与 $\{x:x_i\ge 0,\ i=1,\dots,L\}$ 的交，而每个半空间都是闭集，由有限交闭即闭，故 $B(p,m)$ 闭且有界，由定理~\ref{thm:heineborel} 即紧。向下闭合性直接由 $p\cdot(\lambda x)=\lambda(p\cdot x)\le\lambda m\le m$ 得出。
\end{proof}

把定理~\ref{thm:heineborel} 与命题~\ref{prop:budgetcompact} 合起来读，就能看出本章前面那句“最优解存在”的真实强度：它\emph{不是}“函数总能取到最大”这样的无条件承诺，而是“当偏好可由\emph{连续}效用表示、且预算集\emph{紧}时，取到最大值的点存在”。§1.3 将把三条前提逐一拆开：去掉连续性，最优解可能不存在；去掉紧性（预算集开口），即便效用连续也未必取到。

\begin{figure}[H]
\centering
\begin{tikzpicture}[font=\scriptsize]
  % (a) 闭且有界 → 紧
  \draw[axisline] (0,0) -- (2.7,0) node[below right] {$x_1$};
  \draw[axisline] (0,0) -- (0,2.5) node[above left] {$x_2$};
  \fill[gray!15] (0.3,0.3) rectangle (2.2,1.9);
  \draw[very thick] (0.3,0.3) rectangle (2.2,1.9);
  \fill (1.25,1.1) circle (1.5pt) node[above right=1pt] {$x^{*}$};
  \node[figlabel,align=center] at (1.25,2.28) {(a) 闭且有界 $\Rightarrow$ 紧};
  % (b) 有界不闭
  \begin{scope}[xshift=5.2cm]
    \draw[axisline] (0,0) -- (2.7,0) node[below right] {$x_1$};
    \draw[axisline] (0,0) -- (0,2.5) node[above left] {$x_2$};
    \fill[gray!15] (0.3,0.3) rectangle (2.2,1.9);
    \draw[very thick] (0.3,0.3) -- (2.2,0.3) -- (2.2,1.9);
    \draw[very thick,dashed] (0.3,0.3) -- (0.3,1.9);
    \node[figlabel] at (1.25,0.02) {虚线：未包含的边};
    \fill (1.25,1.1) circle (1.5pt) node[above right=1pt] {$x^{*}$};
    \node[figlabel,align=center] at (1.25,2.28) {(b) 有界不闭 $\Rightarrow$ 不紧};
  \end{scope}
  % (c) 闭但无界
  \begin{scope}[xshift=10.4cm]
    \draw[axisline] (0,0) -- (2.7,0) node[below right] {$x_1$};
    \draw[axisline] (0,0) -- (0,2.5) node[above left] {$x_2$};
    \fill[gray!15] (0.3,0.3) -- (2.4,0.3) -- (2.4,1.9) -- (0.3,1.9) -- cycle;
    \draw[very thick] (0.3,0.3) -- (2.4,0.3) -- (2.4,1.9) -- (0.3,1.9);
    \draw[very thick] (2.4,0.3) -- (2.4,1.9);
    \fill (1.25,1.1) circle (1.5pt) node[above right=1pt] {$x^{*}$};
    \node[figlabel,align=center] at (1.25,2.28) {(c) 闭但无界 $\Rightarrow$ 不紧};
  \end{scope}
  \node[align=center,text width=9.0cm,font=\scriptsize] at (5.4,-1.05) {只有 (a) 是紧集，因而能保证“最优解存在”这条结论。(b) 缺一条边：沿开口方向可以不断改进而永不取到最大；\\
    (c) 向右无界：序列可以逃向无穷远。微观经济学中“偏好无界”“预算集开口导致角点不取”等现象，都是这两类失效在具体模型中的化身。};
\end{tikzpicture}
\caption{紧与不紧：三种集合形态与“最优解存在”的关系}
\end{figure}
\figexplan{三幅图给出同一“有界闭性”的三种组合。只有 (a) 闭且有界，才是紧集，才能保证连续效用下的最大值被取到。(b) 有界但不闭：虚线所示的那条边被排除，偏好若沿该方向改进（哪怕改进幅度越来越小）就会趋近而永不取到，表现为角点解“在边界上不存在”。(c) 闭但无界：可行方向可以一直延伸，$u$ 沿其单调上升，序列逃向无穷远。这两种失效都不是“偏好画错了”，而是集合本身不具备极值定理要求的性质。}

\subsection{序同构与“实数轴太细”的精确含义}

§1.5 末尾我们说过一句直觉：“实数轴太细”。这句话值得精确化，因为它直接决定了可表示性的阶数门槛。

\begin{definition}[全序、序同构、序嵌入]
设 $X$ 上的二元关系 $\prec$ 满足\emph{完全性}（任取 $x\ne y$，或 $x\prec y$ 或 $y\prec x$）与\emph{传递性}，称 $(X,\prec)$ 为\emph{全序}。
若存在双射 $\varphi:X\to Y$ 使 $x\prec y\iff\varphi(x)\prec_{Y}\varphi(y)$，称 $(X,\prec)$ 与 $(Y,\prec_{Y})$ \emph{序同构}。
若只存在单射 $\varphi$ 满足同一条件，称 $\varphi$ 为\emph{序嵌入}，此时称 $X$ 可\emph{嵌入} $Y$。
\end{definition}

\begin{theorem}[Cantor：可数全序可嵌入 \texorpdfstring{$\Q$}{Q}]\label{thm:cantor}
设 $(X,\prec)$ 是\emph{至多可数}的全序。则存在单射 $\varphi:X\to\Q$，使 $x\prec y\iff\varphi(x)<\varphi(y)$。
\end{theorem}

\begin{proof}[构造]
第一步：把 $X$ 枚举为 $x_{1},x_{2},\dots$（允许重复、允许有限），令 $Z_{0}=\emptyset$。
第二步：对 $n\ge 1$，在 $Z_{n-1}$ 的\emph{凸包}内取一个有理数 $q_{n}$，使得当 $x_{n}$ 插入后保持保序——
具体地，若 $x_{n}$ 落在 $Z_{n-1}$ 中相邻两点 $q_{i}<q_{j}$ 之间（端点情形取足够小或足够大的有理数），就令
$q_{n}=(q_{i}+q_{j})/2$。令 $Z_{n}=Z_{n-1}\cup\{q_{n}\}$。
第三步：每个 $q_{n}\in\Q$（$Z_{n-1}$ 有限，中点仍是有理数），且插入位置保证保序，故令 $\varphi(x_{n})=q_{n}$ 即为所需嵌入。
\end{proof}

\begin{corollary}
若 $(X,\prec)$ 是\emph{不可数}的全序（如 $\R$ 自身、或 $\R^{2}$ 上的词典序），则 $X$ 不能嵌入 $\R$ 的任何子集——
特别地，不存在单射 $f:X\to\R$ 保持序。
\end{corollary}

推论的证明只需一步：任何序嵌入都是单射，而 $\R$ 的任意子集都是至多可数的。定理~\ref{thm:cantor} 与这一推论合起来说明：

\emph{“实数轴太细”的确切含义是“可数性”，而不是“密度”。} 稠密性可以靠不断取中点补救（定理~\ref{thm:cantor} 的第二步正是如此），可数性却无法补救。第三章的“可表示必要条件”之所以要求“\emph{至多可数}的序稠密子集”，正是因为效用函数的值域是 $\R$，而 $\R$ 中两两不交的非空开区间族至多可数（每个区间取一个有理数即得单射）。这条线索将在 \S3.1 得到精确形式。

\begin{figure}[H]
\centering
\begin{tikzpicture}[font=\scriptsize]
  % 第一步
  \draw[axisline] (0,0) -- (7.4,0) node[below right] {$\Q$};
  \node[figlabel,anchor=west] at (0.05,1.35) {第 1 步：$x_{1},x_{2}$ 各占其位};
  \fill (1.2,0) circle (1.8pt) node[above=3pt] {$q(x_{1})$};
  \fill (3.4,0) circle (1.8pt) node[above=3pt] {$q(x_{2})$};
  % 第二步
  \node[figlabel,anchor=west] at (0.05,0.72) {第 2 步：插入 $x_{3}$，取中点};
  \fill (1.2,0) circle (1.8pt) node[above=3pt] {};
  \fill (2.3,0) circle (1.8pt) node[above=3pt] {$q(x_{3})$};
  \fill (3.4,0) circle (1.8pt) node[above=3pt] {};
  \draw[very thin,dashed] (1.2,0) -- (3.4,0);
  % 第三步
  \node[figlabel,anchor=west] at (4.30,1.35) {第 3 步：$x_{4}$ 落在 $q(x_{3})$ 之后};
  \draw[axisline] (0,0) -- (7.4,0);
  \fill (1.0,0) circle (1.8pt) node[above=3pt] {};
  \fill (2.3,0) circle (1.8pt) node[above=3pt] {};
  \fill (3.9,0) circle (1.8pt) node[above=3pt] {};
  \fill (4.9,0) circle (1.8pt) node[above=3pt] {$q(x_{4})$};
  \draw[very thin,dashed] (3.9,0) -- (4.9,0);
  \node[align=center,text width=6.6cm,font=\scriptsize] at (3.7,-1.05) {每一步都只在\emph{已有两点之间取中点}，故所有 $q(\cdot)$ 都是有理数；插入位置由序决定，\\
    于是 $\varphi$ 既单射又保序。这正是定理~\ref{thm:cantor} 的构造：稠密性可以逐点补出，可数性则无法补出。};
\end{tikzpicture}
\caption{Cantor 定理的构造：把可数全序逐步嵌入 $\Q$}
\end{figure}
\figexplan{三步构造的核心是“取中点”。第 1 步把 $x_{1},x_{2}$ 放到 $\Q$ 上；第 2 步把 $x_{3}$ 插在 $q(x_{1})$ 与 $q(x_{2})$ 之间，取中点（有理数的中点仍是有理数）；第 3 步把 $x_{4}$ 插在 $q(x_{3})$ 右侧的恰当位置。虚线标出每一步的“待填充区间”。由于每次只在已有两点之间取中点，得到的值全在 $\Q$ 内；而插入位置由 $x$ 的序决定，故保序。图形要点：序的\emph{密度}可以这样一步步补出来，但序的\emph{元素个数}不能——这正是下一节 Afriat 不等式只能在有限数据上求解的原因。}

\subsection{方案空间与可行集}

\begin{definition}[方案空间与可行集]
设所有潜在方案构成集合 $X$，称为\emph{方案空间}（或选择空间、消费空间）。给定参数 $\theta\in\Theta$，决策者实际能从中选择的方案构成\emph{可行集}
\[
B(\theta)\subseteq X.
\]
\end{definition}

可行集\emph{不是}选择结果，而是选择结果可能来自的集合。把二者混同，是初学者最常犯的概念错误之一。

\begin{figure}[H]
\centering
\begin{tikzpicture}[scale=1.0]
  \draw[thick] (0,0) rectangle (6.2,3.6);
  \node[anchor=north west] at (0.15,3.5) {\small $X$（方案空间）};
  \draw[thick,fill=gray!20] (2.5,1.8) ellipse (1.5 and 1.05);
  \node at (2.5,1.8) {\small $B(\theta)$};
  \node[align=left] at (4.7,0.75) {\small $x\notin B(\theta)$：\\ \small 想得但不可行};
\end{tikzpicture}
\caption{方案空间 $X$ 与可行集 $B(\theta)\subseteq X$}
\end{figure}
\figexplan{大矩形为方案空间 $X$（所有潜在方案的集合），其中的椭圆为可行集 $B(\theta)$——给定参数下决策者实际能选的方案。位于椭圆外、矩形内的点 $x$ 属于“想得但不可行”。关键区分：可行集是选择结果可能来自的集合，而不是选择结果本身。}

\begin{example}[预算集]
消费例子中，$L$ 种商品的消费束为 $x=(x_1,\dots,x_L)\in X=\R_+^{L}$。给定价格向量 $p\in\R_{++}^{L}$ 与收入 $m\ge 0$，\emph{预算集}为
\[
B(p,m)=\{\,x\in\R_+^{L}: p\cdot x\le m\,\},
\qquad p\cdot x=\sum_{l=1}^{L}p_lx_l .
\]
当 $L=2$ 时，预算线 $p_1x_1+p_2x_2=m$ 在两轴上的截距分别为 $m/p_1$ 与 $m/p_2$，斜率为 $-p_1/p_2$；预算集是预算线左下方的三角形区域（含边界）。
\end{example}

\begin{figure}[H]
\centering
\begin{tikzpicture}[scale=0.62]
  \draw[->,thick] (0,0) -- (7.4,0) node[below right] {\small $x_1$};
  \draw[->,thick] (0,0) -- (0,5.6) node[above left] {\small $x_2$};
  % 预算集阴影
  \fill[gray!20] (0,0) -- (6,0) -- (0,4.5) -- cycle;
  \draw[very thick] (6,0) -- (0,4.5) node[midway,above right,sloped] {\small $p_1x_1+p_2x_2=m$};
  \node[below] at (6,0) {\small $m/p_1$};
  \node[left] at (0,4.5) {\small $m/p_2$};
  \node at (2.1,1.15) {\small $B(p,m)$};
  \draw[dashed] (2,0) -- (2,1.5) -- (0,1.5);
  \fill (2,1.5) circle (2pt) node[above right] {\small $x$};
\end{tikzpicture}
\caption{二维预算集：截距、斜率与可行域}
\end{figure}
\figexplan{浅灰三角区域（含边界）即预算集 $B(p,m)$，其斜边为预算线 $p_1x_1+p_2x_2=m$，在两轴上的截距分别为 $m/p_1$ 与 $m/p_2$，斜率为 $-p_1/p_2$。虚线给出可行束 $x$ 的坐标读数。价格比决定斜率（机会成本），收入决定截距。}

\begin{proposition}[预算集的基本性质]
设 $p\in\R_{++}^{L}$，$m\ge 0$，则 $B(p,m)$ 是非空、有界、闭且凸的集合。
\end{proposition}

\begin{proof}
\textbf{非空}：$0\in B(p,m)$，因 $p\cdot 0=0\le m$。
\textbf{有界}：对 $x\in B(p,m)$，由 $p_l>0$ 与 $x_l\ge 0$ 得 $x_l\le m/p_l$，故 $\lVert x\rVert\le \sqrt{L}\,\max_l(m/p_l)$。
\textbf{闭}：$B(p,m)$ 是闭集 $\{x:p\cdot x\le m\}$ 与闭集 $\R_+^{L}$ 的交，闭集的交仍闭。
\textbf{凸}：任取 $x,y\in B(p,m)$ 与 $\lambda\in[0,1]$，记 $z=\lambda x+(1-\lambda)y$。则 $z_l=\lambda x_l+(1-\lambda)y_l\ge 0$，且
\[
p\cdot z=\lambda\,(p\cdot x)+(1-\lambda)(p\cdot y)\le \lambda m+(1-\lambda)m=m,
\]
故 $z\in B(p,m)$。
\end{proof}

\begin{remark}
这个命题值得逐字咀嚼：价格严格为正、收入非负、数量非负——三条看似平淡的假设，恰好保证了下一节要用到的“紧且凸”结构。每一条经济假设都在数学上有明确后果。
\end{remark}

\subsection{映射：从方案到数}

\begin{definition}[函数]
设 $X,Y$ 为两个集合。\emph{函数}（映射）$f:X\to Y$ 是一个规则，使每个 $x\in X$ 对应\emph{唯一}的 $y\in Y$，记 $y=f(x)$。$X$ 称\emph{定义域}，$Y$ 称\emph{值域}（陪域）。
\end{definition}

函数有三要素：定义域、值域、对应规则。三者有一个不同，就是不同的函数。

\begin{definition}[单射、满射、双射]
$f:X\to Y$ 称为\emph{单射}，若 $f(x_1)=f(x_2)\Rightarrow x_1=x_2$；称为\emph{满射}，若对每个 $y\in Y$ 都存在 $x\in X$ 使 $f(x)=y$；既是单射又是满射则称\emph{双射}。
\end{definition}

\begin{example}[效用是一种映射]
在经济学中，映射可以把消费束映为一个实数：
\[
u:X\to\R .
\]
例如 Cobb–Douglas 效用 $u(x_1,x_2)=x_1^{\alpha}x_2^{1-\alpha}$（$0<\alpha<1$），或完全替代效用 $u(x_1,x_2)=a x_1+b x_2$。
\end{example}

这里的 $u(x)$ 是对偏好排序的\emph{数值表示}，而不是方案本身。把方案映到实数，目的是把“更好”转化为可运算的序关系——但这一步骤需要行为假设（理性偏好，见第 5 节），并不等于我们观察到了某种内在的心理单位。

\begin{remark}[效用的序数性]
若 $u$ 表示某偏好，则对任意严格递增变换 $g:\R\to\R$，复合函数 $v=g\circ u$ 表示\emph{同一个}偏好，因为 $u(x)\ge u(y)\iff v(x)\ge v(y)$。所以效用只承载\emph{排序}信息：两个消费束谁更好有意义，“好多少”没有意义。这一“序数性”将在第 6 节评价福利比较时再次出现。
\end{remark}

\notehead{本节要点}
\begin{itemize}
  \item 集合语言：$\in,\subseteq,\cup,\cap,\setminus,(\cdot)^{c},\times,\varnothing$；集合相等用双向包含证明；
  \item $\R^{L}$ 中的开/闭、有界、紧、凸——它们分别对应哪些经济含义；
  \item $B(p,m)=\{x\in\R_+^{L}:p\cdot x\le m\}$ 非空、有界、闭、凸（四条逐条可证）；
  \item 函数三要素；效用映射 $u:X\to\R$ 是序数概念。
\end{itemize}

% ==================== 第 3 节：对应（30–50 分钟）====================
\section{从函数到对应}

\subsection{函数与对应的形式差异}

\begin{definition}[对应]
\emph{对应}（correspondence，又称集值映射）$F:X\rightrightarrows Y$ 是一个规则，使每个 $x\in X$ 对应 $Y$ 的一个\emph{子集} $F(x)\subseteq Y$。
\end{definition}

对照函数 $f:X\to Y$（输出一个点），对应 $F:X\rightrightarrows Y$（输出一个集合）。双横箭头 $\rightrightarrows$ 不是装饰：它提醒我们，对应不是“函数写得不够精确”，而是\emph{理论对象本身允许多值}。

\begin{example}
设 $Y=\{a,b,c\}$，$F(x)=\{a,c\}$。这表示 $a$ 与 $c$ 都是可接受的结果。不能擅自写成 $f(x)=a$——那会丢失“并列”“无差异”或“制度性多解”的信息。
\end{example}

函数是对应的\emph{特殊情形}：若对所有 $x$，$F(x)$ 恰含一个元素，才可以把它视为函数。反向不成立：把多值对象硬写成函数，等于在模型里偷偷加了一条“任意决胜规则”。

\begin{figure}[H]
\centering
\begin{tikzpicture}[scale=0.95]
  % 函数
  \draw[thick] (0,0) ellipse (0.8 and 1.3);
  \draw[thick] (3.0,0) ellipse (0.8 and 1.3);
  \node at (0,1.75) {\small $X$};
  \node at (3.0,1.75) {\small $Y$};
  \fill (0,0.6) circle (1.5pt) node[left] {\small $x_1$};
  \fill (0,-0.5) circle (1.5pt) node[left] {\small $x_2$};
  \fill (3.0,0.7) circle (1.5pt) node[right] {\small $f(x_1)$};
  \fill (3.0,-0.4) circle (1.5pt) node[right] {\small $f(x_2)$};
  \draw[->,thick] (0,0.6) -- (2.9,0.68);
  \draw[->,thick] (0,-0.5) -- (2.9,-0.42);
  \node at (1.5,-2.0) {\small 函数 $f:X\to Y$：一个 $x$ 对一个 $y$};
  % 对应
  \begin{scope}[xshift=7.2cm]
    \draw[thick] (0,0) ellipse (0.8 and 1.3);
    \draw[thick] (3.0,0) ellipse (0.8 and 1.3);
    \node at (0,1.75) {\small $X$};
    \node at (3.0,1.75) {\small $Y$};
    \fill (0,0.1) circle (1.5pt) node[left] {\small $x$};
    \fill (3.0,0.8) circle (1.5pt) node[right] {\small $a$};
    \fill (3.0,-0.6) circle (1.5pt) node[right] {\small $c$};
    \draw[->,thick] (0,0.1) -- (2.9,0.78);
    \draw[->,thick] (0,0.1) -- (2.9,-0.58);
    \node at (1.5,-2.0) {\small 对应 $F:X\rightrightarrows Y$：一个 $x$ 对一个\emph{集合}};
  \end{scope}
\end{tikzpicture}
\caption{函数与对应：单值输出与集值输出}
\end{figure}
\figexplan{左图为函数 $f:X\to Y$：左侧定义域中每个点只对应右侧一个点（箭头单发）。右图为对应 $F:X\rightrightarrows Y$：同一个 $x$ 可同时指向多个元素 $a,c$，输出是一个集合。把多值对象硬写成函数，等于在模型里偷偷加入一条任意决胜规则。}

\subsection{选择对应}

\begin{definition}[选择对应]
给定可行集 $B$，\emph{选择对应}定义为
\[
C(B)\subseteq B,
\]
即从可行集中实际被选出的方案所成的集合。若可行集由参数生成，则写作
\[
C:\Theta\rightrightarrows X,\qquad C(\theta)\subseteq B(\theta).
\]
\end{definition}

当多个方案并列最优时，$C(\theta)$ \emph{必须同时包含它们}。

\begin{example}
设 $B=\{a,b\}$，个体对 $a$ 与 $b$ 无差异，则 $C(B)=\{a,b\}$。若只写 $C(B)=a$，就把理论中的非唯一性错误地替换成了一个任意选择规则——模型于是多了不该有的结论。
\end{example}

\subsection{存在性、唯一性与连续性}

三个性质必须分开问：
\begin{itemize}
  \item \textbf{存在性}：$C(\theta)\neq\varnothing$——“有没有最优解？”
  \item \textbf{唯一性}：$C(\theta)$ 恰含一个元素——“最优解是否只有一个？”
  \item \textbf{连续性}：$\theta$ 微调时 $C(\theta)$ 是否平滑变化——这决定了比较静态是否可靠。
\end{itemize}

\begin{theorem}[Weierstrass 极值定理]
若 $u:X\to\R$ 连续且 $B\subseteq\R^{L}$ 是非空紧集，则 $u$ 在 $B$ 上达到最大值与最小值，即存在 $x^{*}\in B$ 使 $u(x^{*})=\max_{x\in B}u(x)$。因此
\[
\argmax_{x\in B}u(x)\neq\varnothing .
\]
\end{theorem}

\begin{example}[开集上最大值可以不存在]
取 $X=\R$，$B=(0,1)$（开区间），$u(x)=x$。$u$ 连续，但 $\sup_{x\in(0,1)}u(x)=1$ 无法达到：对任意 $x\in(0,1)$，总可取 $x'=(x+1)/2>x$。故 $\argmax_{x\in(0,1)}u(x)=\varnothing$。问题正出在 $B$ 不闭，从而不是紧集。
\end{example}

\begin{remark}
连续效用 + 紧致可行集通常只能保证最大值\emph{达到}（存在性），不能保证\emph{唯一}。唯一性需要额外结构：严格凹（或严格拟凹）的目标函数配上凸的可行集，通常带来唯一最优；而线性效用、平坦偏好或非严格凸偏好可能产生\emph{整段}最优解——这正是 $C(\theta)$ 必须是集合而非点的第二个理由（第一个理由是无差异）。
\end{remark}

\subsection{最优值与最优束如何随参数变化}

上一小节保证了最优束\emph{存在}且唯一。本小节追问一个比较静态层面的问题：参数 $\theta$（价格、收入、税基）连续地变动时，$C(\theta)$ 与 $v(\theta)$ 如何随之变动？这不是纯粹的技术问题——经济学的许多“政策冲击效应”正是靠这条通道传导的。

\begin{definition}[对应的连续性（Berge 意义）]
设 $C:\Theta\rightrightarrows X$。称 $C$ 在 $\theta_{0}$ 处\emph{上半连续}（upper hemicontinuous），若对每个满足 $C(\theta_{0})\subseteq O$（$O$ 开）的 $O$，都存在 $\theta_{0}$ 的邻域 $U$，使得 $\theta\in U\Rightarrow C(\theta)\subseteq O$。
称 $C$ 在 $\theta_{0}$ 处\emph{下半连续}（lower hemicontinuous），若对每个满足 $\theta_{0}\in U$（$U$ 开）的 $U$，都存在邻域 $V$，使得 $\theta\in V\Rightarrow C(\theta)\cap U\neq\varnothing$。
\end{definition}

两个条件分别管两件事：上半连续说“$\theta$ 稍一动，选出来的束不会跳出 $C(\theta_{0})$ 附近的紧集”；下半连续说“$C(\theta_{0})$ 里的每一个束，在 $\theta$ 稍动之后都还仍然被选得到”。偏好既不漏（对映）也不溢出（对含），称为\emph{连续}。

\begin{theorem}[最优值函数上半连续]\label{thm:berge}
设 $X$ 是紧集，$u:X\to\R$ 连续，$B(\cdot)$ 关于 $\theta$ 连续，则
\[
v(\theta)=\max_{x\in B(\theta)}u(x)
\]
是 $\theta$ 的上半连续函数。
\end{theorem}

\begin{proof}
记 $x^{*}$ 为 $\theta_{0}$ 处的最优束。给定任意 $\varepsilon>0$：由 $u$ 在紧集 $B(\theta_{0})$ 上一致连续，存在 $\delta>0$，
使当 $x$ 与 $x^{*}$ 的距离小于 $\delta$ 时 $u(x)>u(x^{*})-\varepsilon$；又由 $B(\cdot)$ 在 $\theta_{0}$ 处连续，存在 $\theta_{0}$ 的邻域 $U$，
使 $\theta\in U$ 时 $B(\theta)$ 整体落在 $x^{*}$ 的 $\delta$ 邻域内。于是在 $U$ 上有
\[
v(\theta)\ \ge\ u(x^{*})-\varepsilon\ =\ v(\theta_{0})-\varepsilon .
\]
因 $\varepsilon$ 任意，这正是 $v$ 在 $\theta_{0}$ 处的上半连续性。
\end{proof}

定理~\ref{thm:berge} 的经济学含义是：\emph{最优值可以“跳”，但只能往下跳，且跳跃必然由约束的变动造成，不能由偏好造成}。更精确地说，$v$ 的向下跳跃只能发生在参数使可行集\emph{收缩}的位置（例如税收使预算集突然变小），而不能发生在可行集\emph{扩张}的位置。因此“税收导致消费者剩余下降”这类结论有严格的数学基础。

\begin{remark}[上半连续不是连续]
把 $v$ 说成“连续”是错的。$v$ 一般只是上半连续：在参数空间的某些点上它会\emph{向下跳跃}。典型的例子是配额或从量税——当配额恰好用尽时，价格的微小上升会使可行集突然收缩，最优值随之跳到一个更低的水平。反过来，$v$ 的上包络自动“是”半连续的，这是包络定理的全部内容。
\end{remark}

\begin{figure}[H]
\centering
\begin{tikzpicture}[font=\scriptsize]
  \draw[axisline] (0,0) -- (6.2,0) node[below right] {$\theta$};
  \draw[axisline] (0,0) -- (0,4.0) node[above left] {$v(\theta)$};
  % 上半连续阶梯：向右跳
  \draw[very thick] (0.6,2.6) -- (2.8,2.6);
  \draw[very thick,dashed] (2.8,2.6) -- (2.8,1.4);
  \draw[very thick] (2.8,1.4) -- (5.2,1.4);
  \draw[dashed,thin] (2.8,1.4) -- (2.8,2.6);
  \fill (2.8,1.4) circle (1.5pt);
  \node[figlabel,anchor=west] at (3.0,1.75) {配额用尽：可行集骤然收缩，$v$ 向下跳};
  \node[figlabel,anchor=west] at (0.75,2.85) {约束未变：$v$ 连续};
  \node[figlabel] at (5.2,1.15) {$\theta$};
  \node[align=center,text width=7.2cm,font=\scriptsize] at (3.1,-0.95) {$v$ 上半连续的含义：在跳点 $\theta^{*}$ 处取“跳后”的那一支（$\le v(\theta^{*})$），而不取“跳前”的高支。\\ 图形上就是从 $(2.8,\,1.4)$ 出发向右的实线，而上方那段不是 $v$ 的值。};
\end{tikzpicture}
\caption{最优值函数的向下跳跃：$v(\theta)$ 只上半连续}
\end{figure}
\figexplan{曲线在 $\theta^{*}\le 2.8$ 处连续，然后在 $\theta^{*}$ 处\textbf{向下}跳到较低水平并继续水平。向右的虚线段只用于对照，不是 $v$ 的取值。这正是“半连续”的几何含义：跳点处取较低的一支（$v(\theta^{*})$ 等于右段高度），从而对任意 $\varepsilon>0$，在 $\theta^{*}$ 的邻域内 $v(\theta)>v(\theta^{*})-\varepsilon$ 成立。这种跳跃只能由可行集的收缩（配额用尽、从量税起征）造成，与偏好无关。}

\notehead{本节要点}
\begin{itemize}
  \item $f:X\to Y$ 输出点，$F:X\rightrightarrows Y$ 输出子集；函数是对应的特例；
  \item 选择对应 $C(\theta)\subseteq B(\theta)$ 天生是集值对象；
  \item 存在性靠“连续 + 紧”，唯一性靠“严格凹 + 凸”，二者不可混淆；
  \item 反例：开区间上的线性函数没有最大值。
\end{itemize}

% ==================== 第 4 节：参数与选择（50–70 分钟）====================
\section{参数、可行集与最优选择}

\subsection{参数如何决定可行集}

参数不是模型外的装饰，而是决定机会集合形状的\emph{输入}。一般写作
\[
B:\Theta\rightrightarrows X,\qquad \theta=(p,m,\tau,\omega,\dots),
\]
其中 $\tau$ 可代表税率或补贴，$\omega$ 可代表禀赋。价格与收入决定预算约束；税收可能改变斜率、截距或可行域的分段结构；配额会直接“切掉”一部分集合。

\begin{figure}[H]
\centering
\begin{tikzpicture}[scale=0.72,transform shape]
  % 平移
  \draw[->,thick] (0,0) -- (7.6,0) node[below right] {\small $x_1$};
  \draw[->,thick] (0,0) -- (0,6.2) node[above left] {\small $x_2$};
  \fill[gray!12] (0,0) -- (6.4,0) -- (0,4.8) -- cycle;
  \draw[very thick] (6.4,0) -- (0,4.8);
  \draw[thick,dashed] (4.2,0) -- (0,3.15);
  \node[below] at (6.4,0) {\small $m'/p_1$};
  \node[below] at (4.2,0) {\small $m/p_1$};
  \node[left] at (0,4.8) {\small $m'/p_2$};
  \node[left] at (0,3.15) {\small $m/p_2$};
  \draw[->,thick] (1.7,2.35) -- (2.6,2.95);
  \node at (3.4,-1.4) {\small （a）收入 $m\to m'$：预算线\emph{平移}扩大};
  % 旋转
  \begin{scope}[xshift=9.6cm]
    \draw[->,thick] (0,0) -- (7.6,0) node[below right] {\small $x_1$};
    \draw[->,thick] (0,0) -- (0,6.2) node[above left] {\small $x_2$};
    \fill[gray!12] (0,0) -- (3.4,0) -- (0,4.8) -- cycle;
    \draw[very thick] (6.4,0) -- (0,4.8);
    \draw[thick,dashed] (3.4,0) -- (0,4.8);
    \node[below] at (6.4,0) {\small $m/p_1$};
    \node[below] at (3.4,0) {\small $m/p_1'$};
    \node[left] at (0,4.8) {\small $m/p_2$};
    \draw[->,thick] (4.2,1.5) arc (20:45:2.6);
    \node at (3.4,-1.4) {\small （b）$p_1$ 上升：预算线绕截距\emph{旋转}};
  \end{scope}
\end{tikzpicture}
\caption{参数变化如何改变可行集：平移与旋转}
\end{figure}
\figexplan{（a）收入由 $m$ 增至 $m'$：预算线平行外移（虚线到实线），斜率 $-p_1/p_2$ 不变，两个截距同比例扩大，可行集整体放大。（b）商品 1 价格上升：预算线绕纵轴截距 $m/p_2$ 向内旋转，斜率变陡，横轴截距由 $m/p_1$ 缩至 $m/p_1'$。区分“哪个参数变了、它如何改变可行集”是比较静态的第一步。}

\begin{remark}
收入变化与价格变化对 $B(p,m)$ 的作用在几何上完全不同：$m$ 只改截距不改斜率（机会成本不变）；$p_1/p_2$ 改变斜率（相对价格即机会成本），同时改变一个截距。比较静态的第一步永远是：\emph{哪个参数变了？它怎么改变可行集？}
\end{remark}

\subsection{最优选择的两种输出：\texorpdfstring{$C(\theta)$}{C(θ)} 与 \texorpdfstring{$v(\theta)$}{v(θ)}}

若偏好可由效用函数 $u:X\to\R$ 表示，则选择对应为
\[
\boxed{\,C(\theta)=\argmax\{\,u(x):x\in B(\theta)\,\}\,}
\]

逐字解释：
\begin{itemize}
  \item $C(\theta)$：参数为 $\theta$ 时的\emph{所有}最优方案（一个集合）；
  \item $\argmax$：取的是“使目标达到最大值的\emph{自变量}集合”，而不是最大值本身；
  \item $u(x)$：目标函数；
  \item $x\in B(\theta)$：只能在参数决定的可行集中选择。
\end{itemize}

必须严格区别：
\[
v(\theta)=\max_{x\in B(\theta)}u(x)\in\R
\qquad\text{与}\qquad
C(\theta)=\{\,x\in B(\theta):u(x)=v(\theta)\,\}\subseteq X .
\]
$v(\theta)$ 是\emph{最大值函数}（indirect utility，间接效用），输出一个\emph{数}；$C(\theta)$ 输出达到该数的\emph{所有方案}。最大值回答“最高能达到多少”，argmax 回答“哪些方案达到最高”。

\begin{figure}[H]
\centering
\begin{tikzpicture}[scale=1.0]
  \draw[->,thick] (0,0) -- (7.6,0) node[below right] {\small $x$};
  \draw[->,thick] (0,0) -- (0,4.4) node[above left] {\small $u(x)$};
  % 可行区间
  \draw[very thick] (1.0,0) -- (6.6,0);
  \node[below] at (1.0,0) {\small $a$};
  \node[below] at (6.6,0) {\small $b$};
  \node[below] at (3.8,-0.35) {\small $B=[a,b]$};
  % u 曲线：两个极大，全局峰在左、局部峰在右
  \draw[very thick,domain=1.0:6.6,samples=220,smooth,variable=\t]
    plot ({\t},{3.6*exp(-((\t-2.5)/0.85)^2)+2.7*exp(-((\t-4.9)/0.70)^2)+0.55});
  % 全局最大值 v
  \draw[dashed] (0,4.15) -- (2.5,4.15);
  \node[left] at (0,4.15) {\small $v=\max u$};
  \fill (2.5,4.15) circle (2pt) node[above right=2pt] {\small $x^{*}_1$};
  \draw[dashed] (2.5,4.15) -- (2.5,0);
  \node[below] at (2.5,0) {\small $2.5$};
  % 较低的局部极大：不是 argmax 的成员
  \draw[dashed] (4.9,3.25) -- (4.9,0);
  \fill (4.9,3.25) circle (2pt) node[above right=2pt] {\small $x^{*}_2$};
  \node[below] at (4.9,0) {\small $4.9$};
  \node[align=left] at (5.9,3.7) {\small $\argmax u=\{x^{*}_1\}$（本图唯一）};
\end{tikzpicture}
\caption{最大值 $v$ 是一个数；$\argmax$ 是达到它的自变量集合}
\end{figure}
\figexplan{横轴为方案 $x$，纵轴为效用 $u(x)$，底部的粗线段标出可行区间 $B=[a,b]$。曲线在可行区间内达到的最高高度即 $v=\max u$（一个数，由水平虚线读出）；达到该高度的自变量构成集合 $\argmax u$。图中 $x^{*}_1$ 处达到全局最高，$x^{*}_2$ 只是较低的局部峰值，故 $\argmax u$ 只含 $x^{*}_1$。}

\begin{example}[数值对照]
设 $B=[0,4]$，$u(x)=-(x-2)^2+5$。则 $v=\max u=5$（一个数），$C=\argmax u=\{2\}$（一个集合，本例为单点集）。若改为 $u(x)=\min\{x,2\}$，则 $v=2$ 而 $C=[2,4]$——\emph{整段}都是最优解。
\end{example}

\subsection{为什么 \texorpdfstring{$C(\theta)$}{C(θ)} 必须允许是集合}

第 3 节已给出理由（无差异）。第二个经典来源是\emph{线性效用}：设 $u(x_1,x_2)=x_1+x_2$，$p=(1,1)$，$m=1$。预算线为 $x_1+x_2=1$，效用线与其斜率相同，故预算线上\emph{每一点}都是最优：
\[
C(p,m)=\{\,(x_1,x_2)\in\R_+^2:x_1+x_2=1\,\},
\]
一条线段，而不是一个点。

\begin{figure}[H]
\centering
\begin{tikzpicture}[scale=0.78,transform shape]
  \draw[->,thick] (0,0) -- (7.2,0) node[below right] {\small $x_1$};
  \draw[->,thick] (0,0) -- (0,5.8) node[above left] {\small $x_2$};
  \fill[gray!15] (0,0) -- (5.6,0) -- (0,4.4) -- cycle;
  \draw[very thick] (5.6,0) -- (0,4.4);
  % 线性无差异线（与预算线平行）
  \draw[thick,dashed] (4.0,0) -- (0,3.14);
  \draw[thick,dashed] (2.6,0) -- (0,2.05);
  \node[below] at (5.6,0) {\small $m/p_1$};
  \node[left] at (0,4.4) {\small $m/p_2$};
  \draw[very thick,black] (5.6,0) -- (0,4.4);
  \node[align=left,black] at (5.6,3.6) {\small $C(p,m)$：\\ \small 整条预算线};
\end{tikzpicture}
\caption{线性效用与预算线斜率相同：最优解是\emph{一整段}集合}
\end{figure}
\figexplan{浅灰三角为预算集，其实线斜边为预算线；两条虚线是与预算线平行的线性无差异曲线。由于线性效用 $u=x_1+x_2$ 的边际替代率恒定且恰等于价格比，无差异曲线与预算线斜率相同，预算线上每一点都是最优——选择对应 $C(p,m)$ 是整条预算线（一整段集合），而非单点。}

\subsection{一阶条件、隐函数定理与勒让对偶}

前三小节回答了“最优解存在、唯一、随参数变化”。本小节补上比较静态的数学装置：内点解的一阶条件、隐函数定理保证可微性，以及勒让对偶把“最大化效用”换成“最小化支出”。

\begin{lemma}[内点解的一阶条件]\label{lem:foc}
设 $x^{*}\in\operatorname{int}B(p,m)$ 是 $u$ 的最大点，且 $u$ 可微。则对每一 $i$，
\[
\frac{\partial u}{\partial x_{i}}(x^{*})\ =\ \lambda\,p_{i}\qquad\text{对某个 }\lambda>0,
\]
等价地 $\mathrm{MRS}_{i}=p_{i}/p_{1}$（以商品 1 为计价物）。
\end{lemma}

\begin{proof}
在预算集内部，约束 $\boldsymbol{p}\cdot x=m$ 不起作用，故 $x^{*}$ 是 $u$ 在开集上的局部最大点，由一阶必要条件 $\nabla u(x^{*})=\lambda\nabla(\boldsymbol{p}\cdot x)$，其中拉格朗日乘子 $\lambda$ 由 $x^{*}$ 处于 $m>0$ 且偏好非退化保证为正。等式 $\partial u/\partial x_{i}=\lambda p_{i}$ 两边同除 $\partial u/\partial x_{1}=\lambda p_{1}$ 即得 MRS 表述。
\end{proof}

\begin{theorem}[隐函数定理（经济形式）]
设 $F:\R^{L}\times\R^{2}\to\R$ 连续可微。若 $F(\cdot,\theta^{0})=0$ 在某点 $x^{0}$ 处关于 $x$ 的雅可比可逆，则存在 $\theta^{0}$ 的邻域 $U$ 与 $x^{0}$ 的邻域 $V$，使 $F(x,\theta)=0$ 在 $V\times U$ 内唯一确定一个连续可微函数 $x(\theta)=\bigl(x_{1}(\theta),\dots,x_{L}(\theta)\bigr)$；且
\[
\frac{\partial x}{\partial\theta_{k}}
=\Bigl(\frac{\partial F}{\partial x}\Bigr)^{-1}\frac{\partial F}{\partial\theta_{k}} .
\]
\end{theorem}

把定理用在 $F(x,\theta)=\boldsymbol{p}(\theta)\cdot x-m(\theta)=0$ 上，就得到\emph{需求函数在可行域内连续可微}，从而比较静态可由 $\mathrm{d}x^{*}=\bigl(\partial F/\partial x\bigr)^{-1}\mathrm{d}\theta$ 读出。这是消费者理论中一切“弹性”结论的技术依据。

\begin{theorem}[勒让对偶（支出函数）]\label{thm:leontief}
设 $p\gg 0$、$u$ 连续且\emph{局部非饱和}。定义支出函数
\[
e(p,u)=\min\bigl\{\boldsymbol{p}\cdot x\ :\ x\in\R^{L}_{+},\ u(x)\ge u\bigr\}.
\]
则：$e$ 关于 $p$ 凸、关于 $u$ 递增；$e$ 是唯一支出函数；且当最优解存在唯一时，由包络定理
\[
\frac{\partial e(p,u)}{\partial p_{i}}=x^{*}_{i}\bigl(p,u\bigr),
\qquad
\frac{\partial e(p,u)}{\partial u}=-\lambda\bigl(p,u\bigr)<0 ,
\]
其中 $\lambda$ 是拉格朗日乘子。
\end{theorem}

\begin{remark}[对偶的语言]
原问题问“要达到效用 $u$ 至少要花多少钱”，对偶问题问“要达到支出 $m$ 最多能得多少效用”，即间接效用 $v(p,m)=u\bigl(x^{*}(p,m)\bigr)$。二者的桥梁就是包络定理：\emph{支出对价格的导数恰好是需求，间接效用对价格的导数恰好是 Marshall 需求}。这把“行为”翻译成了“函数”，正是从第一章走向效用理论（第三章）与福利分析（第四章）的接口。
\end{remark}

\begin{figure}[H]
\centering
\begin{tikzpicture}[font=\scriptsize,scale=1.08,transform shape]
  \draw[axisline] (0,0) -- (5.6,0) node[below right] {$m$};
  \draw[axisline] (0,0) -- (0,4.2) node[above left] {$v(p,m)$};
  % 支出函数
  \draw[very thick] (0.5,0.85) .. controls (1.6,1.35) and (2.6,1.75) .. (3.4,2.25)
    node[above right,font=\scriptsize] {$e(p,u)$：凸、递增};
  \fill (2.0,1.15) circle (1.8pt);
  \draw[very thin,dashed] (2.0,0) -- (2.0,1.15) node[below=2pt,font=\scriptsize] {$m$};
  \draw[very thin,dashed] (0,1.15) -- (2.0,1.15) node[left=2pt,font=\scriptsize] {$u$};
  \node[align=left,text width=3.5cm,font=\scriptsize] at (5.0,3.15) {支出函数 $e(p,u)$：\\ 横轴为支出 $m$，纵轴为目标效用 $u$。\\ 它是凸的、单调递增的，\\ 斜率即 $-\partial m/\partial u$（即 $1/\lambda$）};
  \node[align=center,text width=6.4cm,font=\scriptsize] at (2.7,-1.15) {勒让对偶把“最大化 $u(x)$，约束 $\boldsymbol{p}\cdot x\le m$”换成“最小化 $\boldsymbol{p}\cdot x$，约束 $u(x)\ge u$”。\\
    包络定理说：$\partial e/\partial p_{i}=x^{*}_{i}$——最优点积的对价格导数就是需求本身。};
\end{tikzpicture}
\caption{勒让对偶：支出函数及其与最优解的对应}
\end{figure}
\figexplan{支出函数 $e(p,u)$ 以支出 $m$ 为横轴、目标效用 $u$ 为纵轴。图形由左下向右上倾斜且是凸的：要多得一份效用必须多花一笔钱，而且越往上越贵（凸性）。虚线标出任意一点 $(m,u)$ 及其对应的最优束。包络定理的含义是：沿水平方向移动（提高目标效用）所增加的支出，等于原问题中最优消费束的变动量；这正是“需求 = 支出对价格的导数”这条恒等式的几何来源。}

\subsection{比较静态}

\begin{definition}[比较静态]
比较静态（comparative statics）研究参数 $\theta$ 变化如何引起 $B(\theta)$、$C(\theta)$ 或 $v(\theta)$ 的变化。形式上，若 $v$ 可微，则
\[
\mathrm{d}v=\sum_{k}\frac{\partial v}{\partial\theta_k}\,\mathrm{d}\theta_k .
\]
\end{definition}

\begin{example}[收入效应的符号]
在正规品假设下，$\partial x_l^{*}/\partial m>0$：收入上升，需求上升。这是一条\emph{可检验}的比较静态命题——第 6 节会看到，这类命题正是实证研究的典型产出。
\end{example}

\begin{remark}[比较静态 $\neq$ 因果识别]
比较静态给出的是\emph{模型内}结论：“保持其他结构不变，$\theta$ 变，则 $C$ 变”。若现实中价格变化同时携带信息、预期或制度变化，结构解释还需额外假设。博士生应始终追问三问：改变的是哪个原语（primitive）？保持了哪些条件？结论对替代参数化是否稳健？
\end{remark}

\begin{figure}[H]
\centering
\begin{tikzpicture}[scale=0.78,transform shape]
  \draw[->,thick] (0,0) -- (7.4,0) node[below right] {\small $x_1$};
  \draw[->,thick] (0,0) -- (0,5.9) node[above left] {\small $x_2$};
  \fill[gray!12] (0,0) -- (6.2,0) -- (0,4.6) -- cycle;
  \draw[very thick] (6.2,0) -- (0,4.6);
  % 无差异曲线
  \draw[thick,domain=0.35:6.4,samples=100] plot (\x,{1.9/\x});
  \draw[thick,domain=0.6:6.9,samples=100] plot (\x,{3.6/\x});
  \draw[thick,domain=1.0:7.1,samples=100] plot (\x,{6.0/\x});
  % 切点
  \fill (2.35,1.55) circle (2pt);
  \node[fill=white,inner sep=2pt,anchor=west,align=left] at (2.58,1.82) {\small 最优点\\[-1pt]\small $x^{*}=C(p,m)$};
  \draw[dashed] (2.35,1.55) -- (2.35,0) node[below=3pt,fill=white,inner sep=1pt] {\small $x_1^{*}$};
  \draw[dashed] (2.35,1.55) -- (0,1.55) node[left=3pt,fill=white,inner sep=1pt] {\small $x_2^{*}$};
  \node[fill=white,inner sep=2pt,anchor=west] at (5.05,4.35) {\small 无差异曲线};
\end{tikzpicture}
\caption{内点最优：无差异曲线与预算线相切于 $x^{*}$}
\end{figure}
\figexplan{浅灰三角为预算集，其斜边为预算线；三条双曲型无差异曲线由内向外效用递增。内点最优 $x^{*}$ 出现在无差异曲线与预算线相切处（黑点）：切点条件即边际替代率等于价格比。两条虚线给出最优束的两个分量读数 $x_1^{*}$ 与 $x_2^{*}$。}

\notehead{本节要点}
\begin{itemize}
  \item $v(\theta)\in\R$ 是数，$C(\theta)\subseteq X$ 是集合——输出类型不同，不可混写；
  \item $\argmax$ 取“使目标最大的自变量集合”；
  \item 线性效用、平坦偏好 $\Rightarrow$ $C(\theta)$ 可以是一整段；
  \item 比较静态三问：改什么、保什么、是否稳健。
\end{itemize}

% ==================== 第 5 节：纯选择理论（70–85 分钟）====================
\section{纯选择理论：不预设目标函数}

\subsection{公理化起点：只规定选择对应}

抽象选择理论从\emph{可观察的选择事实}出发，只规定三样东西：
\[
X\ (\text{方案空间}),\qquad \mathcal{B}\ (\text{可行集的族}),\qquad C(B)\subseteq B\ (\text{选择对应}),
\]
而\emph{不必先给定}效用函数 $u$。

\subsection{为什么可以不预设 \texorpdfstring{$u$}{u}}

至少有四条理由：
\begin{enumerate}
  \item \textbf{目标函数可能不可观察}。我们能记录“选了什么”，但看不到“为什么”。
  \item \textbf{行为可能由非最优化机制产生}。程序、注意力、社会规范、有限理性都可能生成选择。
  \item \textbf{选择可能循环}。若 $a$ 胜 $b$、$b$ 胜 $c$、$c$ 又胜 $a$，则不存在传递偏好可以理性化这些选择。
  \item \textbf{识别不唯一}。即使存在效用表示，数据也未必能唯一识别它（序数效用的单调变换无穷多）。
\end{enumerate}

因此，“存在某个 $u$ 使 $C(B)=\argmax_{x\in B}u(x)$”是一个\emph{待检验、待证明的理性化命题}，而不是逻辑起点。

\begin{definition}[理性化]
称偏好关系 $\succeq$ \emph{理性化}（rationalize）选择对应 $C$，若对每个可行集 $B$，
\[
C(B)=\{\,x\in B : x\succeq y,\ \forall\, y\in B\,\}.
\]
\end{definition}

\subsection{显示偏好与一致性条件}

显示偏好理论的核心思想是：从\emph{不同预算集}中的实际选择，检验是否存在一个统一偏好关系可以解释全部选择。

\begin{definition}[直接显示偏好]
若 $x$ 在 $y$ 可行时被选择（即 $y\in B$ 而 $x\in C(B)$），称 $x$ \emph{直接显示偏好于} $y$，记 $x\succeq^{R}y$。
\end{definition}

\begin{condition}[WARP，显示偏好弱公理]
若 $x\succeq^{R}y$ 且 $x\neq y$，则不可能有 $y\succeq^{R}x$。换言之：若 $x$ 在 $y$ 买得起时被选中，那么当 $y$ 被选中时，$x$ 必定\emph{买不起}。
\end{condition}

\begin{example}[违反 WARP 的选择]
设 $B_1=\{a,b\}$，观察得 $C(B_1)=\{a\}$（故 $a\succeq^{R}b$）。又设 $B_2=\{a,b,c\}$，若观察得 $b\in C(B_2)$ 而 $a$ 也可行，则 $b\succeq^{R}a$，与 $a\succeq^{R}b$ 矛盾——这组选择不能被任何理性偏好解释。
\end{example}

把 WARP 加强为沿“显示偏好链”的一致性，即得 \textbf{SARP}（显示偏好强公理）：若 $x_1\succeq^{R}x_2\succeq^{R}\cdots\succeq^{R}x_n$ 且 $x_1\neq x_n$，则不能有 $x_n\succeq^{R}x_1$。Afriat 定理表明：满足 SARP 的有限选择数据，必可被某个单调、连续、凹的效用函数理性化——\emph{理性化问题从而有了可检验的答案}。

\begin{figure}[H]
\centering
\begin{tikzpicture}[scale=1.0]
  \fill (0,0) circle (2.2pt) node[below] {\small $a$};
  \fill (3.2,0) circle (2.2pt) node[below] {\small $b$};
  \fill (1.6,2.6) circle (2.2pt) node[above] {\small $c$};
  \draw[->,very thick] (0.35,0.15) -- (2.85,0.15);
  \node at (1.6,-0.5) {\small $a\succ b$};
  \draw[->,very thick] (3.1,0.4) -- (2.0,2.3);
  \node at (3.05,1.6) {\small $b\succ c$};
  \draw[->,very thick] (1.2,2.3) -- (0.15,0.42);
  \node at (0.1,1.6) {\small $c\succ a$};
  \node[align=center] at (1.6,-1.4) {\small 循环偏好 $a\succ b\succ c\succ a$：\\ \small 无传递性，无法被效用函数表示};
\end{tikzpicture}
\caption{循环选择：理性化失败的典型反例}
\end{figure}
\figexplan{三个顶点 $a,b,c$ 与三条箭头构成显示偏好循环：$a\succ b$、$b\succ c$、$c\succ a$。由于传递性不成立，该选择数据无法被任何效用函数表示，也无法通过 SARP 检验——Afriat 定理给出的可检验答案在此失败。这是理性化失败最典型的反例。}

\subsection{理性化如何被检验：Afriat 定理与 GARP}

§1.5 我们只给出“理性化”的定义：存在 $u$ 使选择可由 $u$ 最大化。本小节补上这块拼图的关键一块——\emph{如何用有限观测判断理性化成立}。这是显示偏好理论的核心技术（Afriat 1967），也是把“纯选择理论”从描述性概念变成可操作检验的关键一步。

\begin{definition}[直接显示偏好]
给定有限观测 $\{(p^{t},x^{t})\}_{t=1}^{T}$。若在 $t$ 期 $x^{s}$ 是可负担的（$p^{t}\cdot x^{s}\le p^{t}\cdot x^{t}$），则记 $x^{s}\succeq^{R}x^{t}$，读作“$x^{s}$ 在 $t$ 期显示优于 $x^{t}$”。
\end{definition}

\begin{definition}[GARP（广义显示偏好公理）]
称观测 $\{(p^{t},x^{t})\}_{t=1}^{T}$ 满足 GARP，若对一切 $s\ne t$，
\[
p^{t}\cdot x^{s}\le p^{t}\cdot x^{t}\quad\Longrightarrow\quad
p^{s}\cdot x^{t}\ge p^{s}\cdot x^{s}.
\]
\end{definition}

GARP 的读法：若 $x^{s}$ 在 $t$ 期买得起，那么在 $s$ 期 $x^{t}$ 也必须买得起——否则同一个人先选了 $x^{s}$、后又选了在 $s$ 期买不起 $x^{t}$ 的 $x^{t}$，偏好就出现了不可解释的矛盾。

\begin{remark}[预告：完整结论见 \S2.7.5]
本章只给出这个检验的\emph{思想}：把偏好还原问题变成逐条数据可核验的算术问题。
其完整形式是 \textbf{Afriat 定理}（Afriat 1967）——有限数据可被单调、连续、凹的效用函数
理性化\emph{当且仅当}它满足 GARP；而且一旦可理性化，所需效用可以由一组\emph{线性不等式}
（Afriat 不等式）在有限步内解出，检验的计算量是 $O(n^{3})$（Varian 1982）。
这套算法与它的实证实现（家庭支出数据检验）是 \S2.7 的主题，此处不重复。
\end{remark}

\begin{example}[两期数据的 GARP 检验]\label{ex:garp}
设观测为 $p^{1}=(1,1),\,x^{1}=(1,1)$ 与 $p^{2}=(1,2),\,x^{2}=(1,1)$。
逐对检查（$T=2$，只需查 $s=1,t=2$ 与 $s=2,t=1$）：
\[
p^{2}\cdot x^{1}=1+2=3\le p^{2}\cdot x^{2}=3
\quad\Longrightarrow\quad x^{1}\in B^{2},
\]
故 $x^{1}\succeq^{R}x^{2}$；反向地
\[
p^{1}\cdot x^{2}=1+1=2\le p^{1}\cdot x^{1}=2
\quad\Longrightarrow\quad x^{2}\in B^{1},
\]
故 $x^{2}\succeq^{R}x^{1}$。两条推理首尾相衔，得到 $x^{1}\succeq^{R}x^{2}\succeq^{R}x^{1}$，违反传递性，因而违反 GARP。
由定理~\ref{thm:afriat}，\textbf{不存在任何单调、连续、凹的效用函数能够理性化这两期数据}。
\end{example}

\begin{example}[同一检验通过的情形]\label{ex:garp-ok}
把上例的第二个观测改为 $p^{2}=(1,1),\,x^{2}=(1,1)$。此时 $p^{1}\cdot x^{2}=2=p^{1}\cdot x^{1}$ 仍给出 $x^{2}\succeq^{R}x^{1}$，但反向的 $p^{2}\cdot x^{1}=2>p^{2}\cdot x^{2}=2$ 不成立（严格大于），故 $x^{1}\not\succeq^{R}x^{2}$，不存在环，GARP 通过。数据可被理性化。
\end{example}

例~\ref{ex:garp} 与例~\ref{ex:garp-ok} 的差别仅在价格，但一者被拒、一者被接受——这说明\textbf{价格数据本身就是对偏好的检验}。在没有效用函数、只有价格与数量的观测时，理性化不再是一个不可检验的形而上学命题，而是一个可以逐条数据判真的算术问题。本章到此只给出这个思想的第一次预演；\S2.3--\S2.5 将把它展开为 WARP/SARP/HARP 的完整公理体系，\S2.7 给出 Afriat 定理与可计算的检验算法。

\begin{figure}[H]
\centering
\begin{tikzpicture}[font=\scriptsize]
  % (a) 第 1 期
  \draw[axisline] (0,0) -- (3.0,0) node[below right] {$x_1$};
  \draw[axisline] (0,0) -- (0,2.8) node[above left] {$x_2$};
  \draw[very thick,dashed] (2.0,0) -- (0,2.0);
  \fill (1.0,1.0) circle (2pt) node[above right=2pt] {$x^{1}$};
  \node[figlabel,align=center] at (1.45,2.55) {第 1 期：$p^{1}=(1,1)$\\ $m^{1}=2$，选 $x^{1}$};
  % (b) 第 2 期
  \begin{scope}[xshift=5.4cm]
    \draw[axisline] (0,0) -- (3.0,0) node[below right] {$x_1$};
    \draw[axisline] (0,0) -- (0,2.8) node[above left] {$x_2$};
    \draw[very thick,dashed] (3.0,0) -- (0,1.5);
    \fill (1.0,1.0) circle (2pt) node[above right=2pt] {$x^{2}$};
    \node[figlabel,align=center] at (1.45,2.55) {第 2 期：$p^{2}=(1,2)$\\ $m^{2}=3$，选 $x^{2}$};
  \end{scope}
  \node[figlabel,align=center] at (3.0,-0.50) {$x^{2}\in B^{1}\Rightarrow x^{1}\succeq^{R}x^{2}$};
  \node[figlabel,align=center] at (3.0,-0.95) {$x^{1}\in B^{2}\Rightarrow x^{2}\succeq^{R}x^{1}$};
  \node[align=center,text width=8.6cm,font=\scriptsize] at (3.9,-1.95) {两期所选束相同（都是 $(1,1)$），但两期价格不同：$p^{1}=(1,1)$ 与 $p^{2}=(1,2)$ 都恰好使它用尽预算，\\于是它在对方预算集中都“可负担”，显示偏好双向成立，构成 $2$-环而违反 GARP。};
\end{tikzpicture}
\caption{两期数据构成的 $2$-环：GARP 检验失败的最小实例}
\end{figure}
\figexplan{两幅图分别是两期的预算线（虚线三角）与被选中的束（黑点）。左图第 1 期价格 $(1,1)$、收入 $2$，右图第 2 期价格 $(1,2)$、收入 $3$。两期都恰好选中同一个束 $(1,1)$，且该束在对方的预算集内都“可负担”（支出都正好等于对方收入）。因此两个方向的显示偏好推理都成立，拼成一个 $2$-环，违反 GARP 与传递性，图~\ref{fig:garp-cycle} 给出该组数据不可理性化的结论。}
\label{fig:garp-cycle}

\subsection{把 Afriat 不等式算到底：一个三期数据的完整演算}

上一节把 Afriat 定理的结论预告给了 \S2.7。这里我们把它最核心的计算步骤在三个观测上完整走一遍，使读者看到“检验”究竟是怎样变成算术的。

\begin{definition}[Afriat 不等式（凹效用情形）]\label{def:afriat-ineq}
称效用水平 $(U^{1},\dots,U^{T})\in\R^{T}$ 与斜率 $(\lambda^{1},\dots,\lambda^{T})\in\R_{++}^{T}$ 为观测 $\{(\boldsymbol{p}^{i},x^{i})\}$ 的一组 \textbf{可行解}，若对一切 $i,j$ 都有
\begin{equation}\label{eq:afriat-ineq}
U^{j}\ \le\ U^{i}+\lambda^{i}\,\boldsymbol{p}^{i}\cdot\bigl(x^{j}-x^{i}\bigr).
\end{equation}
\end{definition}

条件~\eqref{eq:afriat-ineq} 的含义是：$u(x)=\min_{i}\bigl\{U^{i}+\lambda^{i}\boldsymbol{p}^{i}\cdot(x-x^{i})\bigr\}$ 是凹的、逐段线性的，而每个 $x^{i}$ 都在 $B^{i}$ 上取到它的上确界。定理~\ref{thm:afriat}（$\S2.7.5$）说：这样的可行解存在\emph{当且仅当} GARP 成立。

\begin{example}[三期数据的手算]\label{ex:afriat3}
设三期观测为
\[
\boldsymbol{p}^{1}=(1,2),\ x^{1}=(2,1);\qquad
\boldsymbol{p}^{2}=(2,1),\ x^{2}=(1,2);\qquad
\boldsymbol{p}^{3}=(1,1),\ x^{3}=(1,1).
\]
\textbf{第一步，逐对检查 GARP。}支出为 $\boldsymbol{p}^{1}\!\cdot x^{1}=\boldsymbol{p}^{2}\!\cdot x^{2}=4$、$\boldsymbol{p}^{3}\!\cdot x^{3}=2$。
可达性：$x^{3}\in B^{1}$（$1\cdot1+2\cdot1=3\le 4$）、$x^{3}\in B^{2}$（$2\cdot1+1\cdot1=3\le 4$）；
而 $x^{2}\notin B^{1}$（$1\cdot1+2\cdot2=5>4$）、$x^{1}\notin B^{2}$（$5>4$）、$x^{1},x^{2}\notin B^{3}$（$3>2$）。
故仅有两条显示偏好：$x^{1}\succeq^{R}x^{3}$ 与 $x^{2}\succeq^{R}x^{3}$，无环，GARP 通过。

\textbf{第二步，解不等式组。}我们试 $U=(U^{1},U^{2},U^{3})=(2,2,0)$。逐对代入条件~\eqref{eq:afriat-ineq}（$p^{i}\cdot(x^{j}-x^{i})$ 的值先算好）：

\begin{center}
\small
\begin{tabular}{ccl}
$(i,j)$ & $\boldsymbol{p}^{i}\cdot(x^{j}-x^{i})$ & 由~\eqref{eq:afriat-ineq} 得到的限制 \\ \midrule
$(1,2)$ & $(1,2)\cdot(-1,\,1)=1$ & $U^{2}\le U^{1}+\lambda^{1}$，即 $2\le 2+\lambda^{1}$，对任意 $\lambda^{1}>0$ 成立 \\
$(1,3)$ & $(1,2)\cdot(-1,\,0)=-1$ & $0\le 2-\lambda^{1}$，即 $\lambda^{1}\le 2$ \\
$(2,1)$ & $(2,1)\cdot(1,\,-1)=1$ & $2\le 2+\lambda^{2}$，对任意 $\lambda^{2}>0$ 成立 \\
$(2,3)$ & $(2,1)\cdot(0,\,-1)=-1$ & $0\le 2-\lambda^{2}$，即 $\lambda^{2}\le 2$ \\
$(3,1)$ & $(1,1)\cdot(1,\,0)=1$ & $2\le 0+\lambda^{3}$，即 $\lambda^{3}\ge 2$ \\
$(3,2)$ & $(1,1)\cdot(0,\,1)=1$ & $2\le 0+\lambda^{3}$，即 $\lambda^{3}\ge 2$ \\
\end{tabular}
\end{center}

取 $\lambda^{1}=\lambda^{2}=1$、$\lambda^{3}=2$，六条限制全部满足（$\lambda^{1},\lambda^{2}\le 2$ 且 $>0$；$\lambda^{3}=2\ge 2$）。因此 $U=(2,2,0)$、$\lambda=(1,1,2)$ 是一组可行解。

\textbf{第三步，写出效用并验证。}$u(x)=\min_{i}\{U^{i}+\lambda^{i}\boldsymbol{p}^{i}\cdot(x-x^{i})\}$ 展开为
\[
u(x_{1},x_{2})=\min\bigl\{\,x_{1}+2x_{2}-2,\ \ 2x_{1}+x_{2}-2,\ \ 2x_{1}+2x_{2}-2\,\bigr\}.
\]
在 $B^{3}$（即 $x_{1}+x_{2}\le 2$）上令 $x_{2}=2-x_{1}$，得 $u=\min\{2-x_{1},\,x_{1},\,2\}$，其最大值为 $1$，恰在 $x_{1}=1$ 处取到，即 $x^{3}=(1,1)$ 被理性化。在 $B^{1}$（$x_{1}+2x_{2}\le 4$）上 $u$ 的最大值为 $2$，在从 $x^{1}=(2,1)$ 到 $(4,0)$ 的整段线段上都取到，故 $x^{1}$ 是理性化选择；$x^{2}$ 同理。
\end{example}

例~\ref{ex:afriat3} 显示了三件事。其一，检验是\emph{有限}的：$T=3$ 只涉及 $9$ 个不等式与 $6$ 个变量，人手可算。其二，效用函数由可行解\emph{构造}出来，不是猜出来的。其三，理性化允许最优解\emph{不唯一}——$B^{1}$ 上整段线段都取到最大值 $2$，而 $x^{1}$ 只是其中一个点。这正是“$C(\theta)$ 必须是集合而非点”的第一个具体证据。

\begin{figure}[H]
\centering
\begin{tikzpicture}[font=\scriptsize]
  % 三期预算线与被选束
  \draw[axisline] (0,0) -- (5.4,0) node[below right] {$x_1$};
  \draw[axisline] (0,0) -- (0,4.0) node[above left] {$x_2$};
  \draw[thick,dashed] (4.0,0) -- (0,2.0) node[above right,font=\scriptsize] {$B^{1}$};
  \draw[thick,dashed] (0,4.0) -- (4.0,0) node[below right,font=\scriptsize] {$B^{2}$};
  \draw[thick,dashed] (2.0,0) -- (0,2.0) node[above left,font=\scriptsize] {$B^{3}$};
  \fill (2.0,1.0) circle (2pt) node[above right=1pt] {$x^{1}$};
  \fill (1.0,2.0) circle (2pt) node[above left=1pt] {$x^{2}$};
  \fill (1.0,1.0) circle (2pt) node[below left=1pt] {$x^{3}$};
  \node[align=center,text width=7.2cm,font=\scriptsize] at (2.7,-1.45) {三期预算线与被选中的束（黑点）。显示偏好只有两条：$B^{1}$ 与 $B^{2}$ 都装得下 $x^{3}$，而 $x^{1}$ 与 $x^{2}$ 互相装不下。\\构造出的 $u$ 在 $B^{1}$ 上于 $(2,1)$ 到 $(4,0)$ 的整段线段取到最大值，故最优解不唯一。};
\end{tikzpicture}
\caption{三期数据的 GARP 检验与 Afriat 不等式的解}
\end{figure}
\figexplan{三根虚线分别是三期预算线 $B^{1}$（斜率 $-1/2$）、$B^{2}$（斜率 $-2$）、$B^{3}$（斜率 $-1$），黑点是被选中的束。$x^{3}=(1,1)$ 落在 $B^{1}$ 与 $B^{2}$ 内却不落在 $B^{3}$ 内，$x^{1}$ 与 $x^{2}$ 互不在对方的预算集内——这解释了为什么只有两条显示偏好且无环。图形要点：Afriat 不等式给出的 $u$ 在 $B^{1}$ 上有一段“高原”，最优解因此不唯一；这正是选择对应必须是集合的技术理由。}

\subsection{不预设目标函数 ≠ 否认行为规律}

不预设 $u$，只是\emph{把解释任务后移}：先让数据说话，再问何种偏好（如果存在）能解释它。这是理论经济学的方法论自觉——模型应当对“可观察的不一致”保持敏感，而不是用数学形式掩盖它。

\notehead{本节要点}
\begin{itemize}
  \item 纯选择理论只给 $C(B)\subseteq B$，不预设 $u$；
  \item “$C$ 可被效用理性化”是定理/检验的结论，不是公理；
  \item 直接显示偏好 $\succeq^{R}$、WARP、SARP 构成可检验的一致性阶梯；
  \item 循环偏好 $a\succ b\succ c\succ a$ 是理性化失败的最小反例。
\end{itemize}

% ==================== 第 6 节：实证与规范（85–102 分钟）====================
\section{实证分析与规范分析：两种不同的问题}

经济学研究常被混作一谈的两类问题，其实属于不同的逻辑类别：\emph{实证}（positive）研究“是什么”，\emph{规范}（normative）研究“应当怎样”。把二者分开，是博士生必须具备的方法论基本功。

\subsection{实证分析：回答“是什么”}

\begin{definition}[实证命题]
\emph{实证命题}是关于世界\emph{是}什么样、以及“若……则……”会怎样变化的可检验陈述。它回答三类问题：是什么（描述）、为什么（机制）、会怎样（预测）。
\end{definition}

实证命题的三个特征：
\begin{enumerate}
  \item \textbf{概念清晰}：每个术语都有可操作的定义（如“正常品”定义为 $\partial x_l^{*}/\partial m>0$）；
  \item \textbf{机制明确}：因果链条在模型内部可追踪（$\theta\to B(\theta)\to C(\theta)$）；
  \item \textbf{预测可检验}：原则上存在数据可以\emph{推翻}它。
\end{enumerate}

\begin{example}[比较静态是典型实证命题]
“收入增加会提高对正常品的需求”——在保持偏好、价格不变且满足假设 $A$ 时，参数 $m$ 的上升导致选择集合发生特定变化：
\[
m'\ge m \;\Longrightarrow\; x_l^{*}(p,m')\ge x_l^{*}(p,m)\quad(\text{正常品}).
\]
又如：“税率 $\tau$ 上升会降低劳动供给”——这是关于 $C(\theta)$ 对 $\tau$ 的比较静态，原则上可用微观数据检验。
\end{example}

实证结论的规范写法是带\emph{条件}的：
\[
\text{在保持 }Z\text{ 不变且满足假设 }A\text{ 时，参数 }p\text{ 的变化导致选择发生某种变化。}
\]
丢掉条件，实证命题就从“可检验”退化为“口号”。

\subsection{规范分析：回答“应当怎样”}

\begin{definition}[规范命题]
\emph{规范命题}是关于“什么应当发生”“哪个结果更好”的陈述。它必须援引某个\emph{价值标准}（福利准则），而价值标准不能从事实中自动推出。
\end{definition}

最常用的价值标准有：
\begin{itemize}
  \item \textbf{帕累托标准}：状态 $x$ 帕累托优于 $y$，若无人变差且至少一人变好（严格写法见下）。帕累托改进得到的是\emph{效率}判断，不涉及分配；
  \item \textbf{功利主义社会福利函数}：
        \[
        W(u_1,\dots,u_n)=\sum_{i=1}^{n}\alpha_i u_i,\qquad \alpha_i\ge 0,
        \]
        把社会评价写成个体效用的加权和；
  \item \textbf{罗尔斯（max–min）标准}：$W=\min_i u_i$，社会评价取决于最不利者。
\end{itemize}

\begin{remark}[权重不是事实]
功利主义中的权重 $\alpha_i$ \emph{不是}由个体选择事实自动推出的，它体现的是评价者的价值立场（谁更重要）。把某个 $\alpha$ 当作“客观的”，正是把规范伪装成实证。同理，效率提高不意味公平改善；观察到某政策“有效”，也不等于该政策“正当”。
\end{remark}

\begin{definition}[帕累托最优]
配置 $x$ 称为\emph{帕累托最优}（Pareto optimal），若不存在其他可行配置 $x'$ 使所有人 $i$ 满足 $u_i(x'_i)\ge u_i(x_i)$ 且至少一个 $j$ 满足 $u_j(x'_j)>u_j(x_j)$。
\end{definition}

\begin{example}[效率与公平可以背离]
把所有资源给一个人而其余人为零，可以是帕累托最优（无法在不损害该人的前提下改善他人），但几乎无人认为它“应当”如此。帕累托标准对分配沉默，这正是它作为\emph{最低限度}价值标准的原因与局限。
\end{example}

\subsection{三层分离}

政策评价应分三层进行：

\begin{figure}[H]
\centering
\begin{tikzpicture}[
  box/.style={draw,thick,rounded corners=2pt,minimum width=3.3cm,minimum height=1.05cm,align=center,fill=gray!8},
  lab/.style={font=\small,align=left,text width=4.6cm}
]
  \node[box] (a) at (0,0) {\textbf{事实层}\\ \small 政策改变了什么？};
  \node[box] (b) at (4.6,0) {\textbf{机制层}\\ \small 经何种约束/激励/均衡传导？};
  \node[box] (c) at (9.2,0) {\textbf{价值层}\\ \small 依何种福利标准评价？};
  \draw[->,very thick] (a) -- (b);
  \draw[->,very thick] (b) -- (c);
  \node[lab] at (0,-1.35) {\small 例：最低工资使低技能就业下降 2\%};
  \node[lab] at (4.6,-1.35) {\small 例：劳动需求弹性与覆盖范围决定效应大小};
  \node[lab] at (9.2,-1.35) {\small 例：用 $W=\sum\alpha_i u_i$ 还是罗尔斯标准？};
\end{tikzpicture}
\caption{政策评价的三层分离}
\end{figure}
\figexplan{三个方框依次是政策评价的三层：事实层（政策改变了什么）、机制层（经何种约束、激励或均衡传导）、价值层（依何种福利标准评价）。箭头表示论证推进方向，每层下方给出最低工资的例子。关键纪律：不得把价值层伪装成事实层——由“提高了效率”推不出“应当实行”。}

\begin{remark}
把第三层（价值）伪装成第一层（事实），是经济学论证中最常见的越界。例如“最低工资提高了效率，所以应该实行最低工资”——前半句即使有证据，后半句仍需额外的价值前提才能成立。
\end{remark}

\subsection{案例演练：把一个政策问题拆成正反两半}

\begin{center}
\small
\begin{tabular}{p{2.6cm}p{5.6cm}p{5.6cm}}
\toprule
政策 & 实证命题（是什么） & 规范判断（应当怎样） \\
\midrule
最低工资 & 提高最低工资降低低技能工人的就业量（弹性可估计） & 应当/不应当实行最低工资（需福利标准与分配权重） \\
房产税 & 开征房产税降低投资性购房需求、改变地方财政收入结构 & 应当/不应当开征（效率与公平的权衡是价值问题） \\
消费券/补贴 & 消费券在短期内提高边际消费倾向高的群体的消费 & 应当/不应当发放（跨期效率、财政成本与目标人群是价值判断） \\
\bottomrule
\end{tabular}
\end{center}

注意左栏每一条原则上都可被数据\emph{证伪}；右栏每一条都必须明示其价值标准。两栏不可互相替代。

\begin{example}[最低工资：从实证估计到规范判断]
以美国青少年低技能劳动力市场为例。实证研究关心：最低工资提高 10\% 会导致就业下降多少？早期时间序列估计给出就业弹性约为 $-0.1$ 到 $-0.3$（提高工资显著减少就业）；而 Card \& Krueger（1994）利用新泽西与宾州快餐店的自然实验发现，新泽西最低工资从 \$4.25 提高到 \$5.05 后，就业变化并不显著为负，甚至略升。后续元分析（Neumark \& Wascher, 2007）指出：对青少年就业的弹性中位数约为 $-0.1$，且对低技能群体的负面影响更明显。

这些估计回答的是事实层与机制层。但“是否应当提高最低工资”仍取决于价值判断：若采用功利主义并赋予低收入者较高权重 $\alpha_i$，就业的轻微损失可能被低收入者工资收益抵消；若采用帕累托标准，则任何因失业受损的工人都构成反对理由，除非对他们进行补偿。实证结果缩小了规范争论的范围，却不能单独决定政策。
\end{example}

\subsection{实证与规范的互补}

实证与规范不是对手，而是完整研究的两个半部：
\[
\underbrace{\text{实证：预测政策后果}}_{\text{没有它，规范是空中楼阁}}
\;+\;
\underbrace{\text{规范：评价后果好坏}}_{\text{没有它，实证无政策含义}}
\;\Longrightarrow\; \text{完整的政策分析}.
\]
实证错误会让规范结论建立在沙上；规范模糊会让实证研究失去方向。博士生的写作规范是：每一段都先说清楚自己在哪一层。

\notehead{本节要点}
\begin{itemize}
  \item 实证：是什么/为什么/会怎样；三特征——概念清晰、机制明确、可检验；
  \item 规范：应当怎样；必须援引价值标准（帕累托、功利主义 $W=\sum\alpha_i u_i$、罗尔斯）；
  \item 权重 $\alpha_i$ 是价值立场而非事实；效率 $\neq$ 公平；
  \item 三层分离：事实层、机制层、价值层；越界是最常见错误。
\end{itemize}

% ==================== 第 7 节：模型价值（102–115 分钟）====================
\section{模型的价值与适用边界}

一个好模型不是“更复杂”或“更现实”的模型。本节给出六个评价维度，并说明它们之间的张力。

\subsection{六维评价框架}

\begin{center}
\small
\begin{tabular}{p{2.4cm}p{3.9cm}p{4.0cm}p{3.2cm}}
\toprule
维度 & 核心问题 & 证据 & 主要风险 \\
\midrule
Useful 有用 & 能回答重要问题吗？ & 解释、政策用途、反事实 & 问题选择过窄 \\
Accurate 准确 & 预测与数据一致吗？ & 样本外预测、实验与准实验 & 过拟合 \\
Broad 广泛 & 能迁移到其他场景吗？ & 消费、生产、劳动等应用 & 过宽而不可证伪 \\
Compact 简洁 & 以少量结构产生清晰推论？ & 假设–结论压缩比 & 忽略关键机制 \\
Identifiable 可识别 & 参数能被数据区分吗？ & 识别策略、敏感性分析 & 依赖不可观测假设 \\
Falsifiable 可证伪 & 什么结果会推翻它？ & 明确反事实与拒绝条件 & 模型永远可解释 \\
\bottomrule
\end{tabular}
\end{center}

\begin{figure}[H]
\centering
\begin{tikzpicture}[font=\scriptsize]
  % 图例
  \fill[gray!25] (0,3.34) rectangle (0.22,3.56);
  \node[anchor=west] at (0.30,3.45) {模型 A：简洁但问题过窄};
  \fill[gray!25] (2.62,3.34) rectangle (2.84,3.56);
  \node[anchor=west] at (2.92,3.45) {模型 B：广泛但不可证伪};
  \fill[gray!25] (2.40,{2.74+0.06}) rectangle ({2.40+1.30},2.80);
  \fill[gray!55] (2.40,{2.74-0.10}) rectangle ({2.40+2.55},2.70);
  % 六个维度
  \foreach \y/\nm/\a/\b in {2.70/Useful/1.30/2.55, 2.25/Accurate/2.10/1.55, 1.80/Broad/0.65/2.95,
                            1.35/Compact/2.90/1.10, 0.90/Identifiable/1.35/2.35, 0.45/Falsifiable/0.95/0.70}{
    \node[anchor=east] at (2.30,\y) {\nm};
    \fill[gray!25] (2.40,{\y+0.06}) rectangle ({2.40+\a},{\y+0.20});
    \fill[gray!55] (2.40,{\y-0.10}) rectangle ({2.40+\b},{\y+0.04});
  }
  \node[align=center,font=\scriptsize] at (3.60,-0.28) {条长表示该维度的相对得分。\\两个模型各擅一端：六维之间存在张力，\\“更好的模型”往往只是换了取舍。};
\end{tikzpicture}
\caption{六维评价框架：简洁性与广泛性之间的张力}
\end{figure}
\figexplan{每一行是一个评价维度，浅灰条为模型 A、深灰条为模型 B。模型 A 在 Compact（简洁）上得分高、在 Broad（广泛）上得分低；模型 B 恰好相反。可识别性与可证伪性亦呈此消彼长。图形要传达的不是“哪个模型更好”，而是六维之间不可同时最大化的取舍结构。}

\subsection{逐维展开}

\notehead{7.2.1 Useful：有用而非万能}
模型的有用性取决于\emph{问题}。选择模型能从选择恢复偏好、组织制度比较，并与个体自主性相容；但人际效用比较困难——不能把个体内部的序数排序直接加总为社会福利（这正是第 6 节 $W$ 需要外生权重的根源）。

\notehead{7.2.2 Accurate：准确必须指向可检验含义}
一个模型不是因为数学精致就准确。应区分四种“准确”：样本内拟合、样本外预测、结构参数识别、反事实有效性。比较静态只有在参数具有清晰制度含义、且其他机制被控制时，才有经验内容。

\notehead{7.2.3 Broad 与 Compact 的张力}
越宽泛的模型越可能迁移，但越难产生锐利预测；越紧凑的模型越便于推导，却可能忽略注意力、身份、规范与非传递性。简化本身不是缺陷，关键是明确回答：\emph{为了哪个问题，牺牲了哪些现实因素？}

\notehead{7.2.4 Identifiable：参数能否被数据区分}
若两组截然不同的参数 $(\theta_1,\theta_2)$ 产生完全相同的选择数据，则该参数不可识别。识别失败时，一切“结构解释”都是空中楼阁。补救办法：补充数据维度、施加可辩护的约束、或改问可识别的问题。

\notehead{7.2.5 Falsifiable：必须说清什么会推翻它}
一个“无论看到什么都能解释”的模型没有信息量。写模型时就应写下拒绝条件：若数据呈现 $X$，则放弃该模型。可证伪性是把实证研究与讲故事区分开的底线。

\subsection{简化假设、反例与边界}

经典选择框架的最小假设可表述为：
\begin{enumerate}
  \item 选择来自某个可行集；
  \item 被偏好的方案更可能被选择；
  \item 任意两个潜在方案可以比较（完全性）；
  \item 偏好具有传递性。
\end{enumerate}

每一条假设都应对应一个\emph{结论依赖点}和一个\emph{最小反例}：
\begin{itemize}
  \item 删除\emph{完全性} $\Rightarrow$ 某些方案无法比较，$C(B)$ 可能无法定义；
  \item 删除\emph{传递性} $\Rightarrow$ 可能出现 $a\succ b\succ c\succ a$ 的循环，无法被效用表示；
  \item 删除\emph{紧致性} $\Rightarrow$ 最大值可能不存在（开区间上的线性函数）；
  \item 删除\emph{严格凹性} $\Rightarrow$ 最优方案可能不唯一（线性效用）。
\end{itemize}

\notehead{边界声明的写法}
好的论文在模型之后紧跟一段“边界声明”：本模型在哪些条件下成立、哪些机制被有意忽略、哪些结论对哪些假设敏感。这不是示弱，而是把模型的\emph{适用范围}写清楚——它恰恰是可信度的来源。

\notehead{本节要点}
\begin{itemize}
  \item 六维评价：Useful、Accurate、Broad、Compact、Identifiable、Falsifiable；
  \item Broad 与 Compact 存在张力；Accurate 必须指向可检验含义；
  \item 每条假设对应一个结论依赖点 + 一个最小反例；
  \item 边界声明是模型可信度的一部分。
\end{itemize}

% ==================== 第 8 节：总结、测验与作业（115–120 分钟）====================
\section{总结、课堂测验与作业}

\subsection{一条逻辑链}

本章的核心不是“把一切行为都写成最大化”，而是理解一连串\emph{逻辑转换的条件}：
\[
\underbrace{X\supseteq B(\theta)}_{\text{从方案到可行集}}
\;\longrightarrow\;
\underbrace{C(\theta)\subseteq B(\theta)}_{\text{从可行集到选择}}
\;\longrightarrow\;
\underbrace{\succeq^{R}}_{\text{从选择到偏好}}
\;\longrightarrow\;
\underbrace{u:X\to\R}_{\text{从偏好到效用}}
\;\longrightarrow\;
\underbrace{C(\theta)=\argmax_{x\in B(\theta)}u(x)}_{\text{从效用表达到优化}},
\]
压缩成一句话：
\[
\text{choice data}\;\Rightarrow\;\text{preference representation}\;\Rightarrow\;\text{optimization}\;\Rightarrow\;\text{testable implications}.
\]

\begin{figure}[H]
\centering
\begin{tikzpicture}[font=\small,
  box/.style={draw,thick,rounded corners=2pt,align=center,inner sep=4pt,fill=gray!12,minimum width=2.0cm,minimum height=1.0cm},
  arr/.style={-{Stealth[length=2.4mm]},thick}]
  \node[box] (b1) at (0,0) {$X\supseteq B(\theta)$};
  \node[box] (b2) at (2.5,0) {$C(\theta)\subseteq B(\theta)$};
  \node[box] (b3) at (5.0,0) {$\succeq^{R}$};
  \node[box] (b4) at (7.5,0) {$u:X\to\R$};
  \node[box] (b5) at (10.0,0) {$C(\theta)=\argmax u$};
  \draw[arr] (1.05,0) -- (1.45,0);
  \draw[arr] (3.55,0) -- (3.95,0);
  \draw[arr] (6.05,0) -- (6.45,0);
  \draw[arr] (8.55,0) -- (8.95,0);
  \node[font=\scriptsize,above=3pt] at (1.25,0) {给出偏好};
  \node[font=\scriptsize,above=3pt] at (3.75,0) {需要理性};
  \node[font=\scriptsize,above=3pt] at (6.25,0) {需效用存在};
  \node[font=\scriptsize,above=3pt] at (8.75,0) {需单调};
  \node[align=center,font=\scriptsize] at (0,-1.05) {集合包含\\（无需偏好）};
  \node[align=center,font=\scriptsize] at (2.5,-1.05) {单值化假设\\（几何直觉）};
  \node[align=center,font=\scriptsize] at (5.0,-1.05) {传递性\\（Varian 修正）};
  \node[align=center,font=\scriptsize] at (7.5,-1.05) {完备性与连续性\\（可数序稠密）};
  \node[align=center,font=\scriptsize] at (10.0,-1.05) {单调性\\（强单调）};
  \node[font=\scriptsize] at (5.0,-2.05) {四个箭头各自要求一个额外条件；任一环失效，整条链即在相应处断裂。};
\end{tikzpicture}
\caption{一条逻辑链：从方案空间到优化表述，及其每一步的前提}
\end{figure}
\figexplan{五个方框自左至右是本章的形式化主线：方案空间含可行集、可行集映到选择、选择反推偏好、偏好由效用表示、效用写成最大化表述。箭头上方标出该步所需的前提，方框下方标出该步真正用到的结构条件。图形要点：链条不是自动成立的，每推进一格都要额外付出一个假设。}

\subsection{三个必须记住的反例}
\begin{enumerate}
  \item \textbf{非唯一最优}：$u(x_1,x_2)=x_1+x_2$，$p=(1,1)$，$m=1$ 时 $C(p,m)$ 是整条预算线——故选择对应\emph{必须}是集值对象；
  \item \textbf{最大值不存在}：$B=(0,1)$、$u(x)=x$ 时 $\argmax=\varnothing$——存在性依赖紧致性；
  \item \textbf{循环偏好}：$a\succ b\succ c\succ a$——传递性缺失时无法被效用理性化。
\end{enumerate}

\subsection{课堂测验（5 分钟）}
\begin{enumerate}
  \item 若两个消费束都达到最大效用，为什么不能把选择对应写成函数？\\
        \emph{要点：}函数要求唯一输出；并列最优意味着输出含两个元素，写成函数等于偷偷加入任意决胜规则，丢失无差异信息。
  \item $v(\theta)$ 与 $C(\theta)$ 的输出类型分别是什么？\\
        \emph{要点：}$v(\theta)=\max u\in\R$ 是数；$C(\theta)=\argmax u\subseteq X$ 是集合。
  \item 预算线移动时，变化的是偏好还是可行集？\\
        \emph{要点：}可行集 $B(\theta)$。偏好 $\succeq$ 是行为原语，不随参数变化；这是比较静态的前提。
  \item “税收提高效率”为什么还不足以推出“税收应当提高”？\\
        \emph{要点：}前者是实证命题（事实层），后者是规范判断（价值层），需要额外福利标准才能连接。
  \item 一个模型越现实是否一定越有价值？\\
        \emph{要点：}否。Broad 与 Compact 存在张力；价值取决于六维综合表现与所回答的问题。
\end{enumerate}

\subsection{课后作业}
\begin{enumerate}
  \item 构造一个最小例子，使 $\argmax$ 恰好含有\emph{两个}元素，并画出图形；
  \item 给出一个没有最大值的可行集与连续目标函数（不得照抄课堂反例），并说明缺少的是哪一条性质；
  \item 用预算集解释收入效应与替代效应的逻辑区别（提示：先分解“哪个参数变了”）；
  \item 选择一个你关心的政策模型，分别写出它的：正向命题、规范标准、识别风险与可证伪条件。
\end{enumerate}

\subsection{结语}

最终评价一个模型时，应同时追问：它是否\emph{有用}、是否\emph{准确}、是否足够\emph{广泛}、是否保持\emph{简洁}、是否\emph{可识别}、是否\emph{可证伪}，以及它的规范结论究竟依赖了哪些价值前提。把这六问变成肌肉记忆，是本课最重要的产出。

% 现有内容的证明补充：并入本章正文层级，不单列附录章节
\subsection{本节证明补充：序列、收敛与需求的局部结构}

\subsection{序列语言与经济学中的极限}
设 $x^n\in X\subseteq\mathbb R^L$ 是一列消费束。若对任意 $\varepsilon>0$，存在 $N$ 使得 $n\ge N$ 时
$\|x^n-x^*\|<\varepsilon$，则称 $x^n\to x^*$。这一表述把“足够接近”变成可检验的量化命题。
在经济学中，序列不仅是分析工具，也代表逐步调整的市场过程：价格、收入、禀赋或政策参数都可沿序列变化。

\begin{definition}[上极限与下极限]
对有界实数列 $(a_n)$，定义
\[
\limsup a_n=\lim_{N\to\infty}\sup_{n\ge N}a_n,\qquad
\liminf a_n=\lim_{N\to\infty}\inf_{n\ge N}a_n.
\]
当且仅当二者相等时，序列收敛。
\end{definition}

\begin{proposition}[连续函数保持序列收敛]
若 $f:X\to\mathbb R$ 在 $x^*$ 连续且 $x^n\to x^*$，则 $f(x^n)\to f(x^*)$。
\end{proposition}
\begin{proof}
给定 $\varepsilon>0$，由连续性存在 $\delta>0$，使 $\|x-x^*\|<\delta$ 时
$|f(x)-f(x^*)|<\varepsilon$。由收敛性存在 $N$，使 $n\ge N$ 时 $\|x^n-x^*\|<\delta$，代入即得。
\end{proof}

\begin{figure}[H]\centering
\begin{tikzpicture}[x=1.0cm,y=0.8cm,font=\scriptsize]
\draw[->] (0,0)--(8,0) node[right] {$n$};\draw[->] (0,0)--(0,4) node[above] {$x^n$};
\draw[dashed] (0,1.5)--(8,1.5) node[right] {$x^*$};
\foreach \n/\y in {0.6/3.2,1.2/2.6,1.8/2.0,2.4/2.8,3.0/1.8,3.6/2.2,4.2/1.7,4.8/1.65,5.4/1.55,6.0/1.52,6.6/1.5,7.2/1.5}{\fill (\n,\y) circle (1.4pt);}
\node[align=center,fill=white] at (4.8,3.25) {震荡幅度逐步缩小\\但未必单调};
\end{tikzpicture}
\caption{序列收敛：调整过程可以非单调}
\end{figure}
\figexplan{横轴是调整次数 $n$，纵轴是某个标量状态变量。点列可上下震荡，但只要最终落入任意给定的窄带，就收敛于虚线所示的 $x^*$。}

\subsection{预算集的边界、内点与切锥}
对价格 $p\gg0$ 和财富 $w>0$，预算集为 $B(p,w)=\{x\ge0:p\cdot x\le w\}$。其边界
$\partial B=\{x:p\cdot x=w\}$ 是预算耗尽集合；内点 $\{x:p\cdot x<w\}$ 则保留小幅调整空间。
若 $x^*$ 位于边界，允许方向 $d$ 必须满足 $p\cdot d\le0$（局部一阶近似），这给出可行方向锥。

\begin{figure}[H]\centering
\begin{tikzpicture}[x=0.9cm,y=0.9cm,font=\scriptsize]
\fill[gray!12] (0,0)--(0,4)--(4,0)--cycle;
\draw[->] (0,0)--(4.7,0) node[right] {$x_1$};\draw[->] (0,0)--(0,4.7) node[above] {$x_2$};
\draw[thick] (0,4)--(4,0) node[pos=.2,above left] {$p\cdot x=w$};
\fill (2.2,1.8) circle (2pt) node[below right] {$x^*$};
\draw[->,dashed] (2.2,1.8)--(3.0,1.4) node[right] {$d_1$};
\draw[->,dashed] (2.2,1.8)--(1.7,2.7) node[left] {$d_2$};
\node[align=center,fill=white] at (1.0,3.6) {可行集\\$p\cdot x\le w$};
\end{tikzpicture}
\caption{预算边界与局部可行方向}
\end{figure}
\figexplan{灰色三角形是预算集，粗线是边界。点 $x^*$ 处的虚线箭头表示局部可行方向；若方向指向灰色区域外，则不能作为一阶可行变动。}

\begin{theorem}[两商品内点需求的一阶条件]
若 $u$ 可微、严格拟凹，且最优束 $x^*\gg0$，则存在 $\lambda>0$ 使
\[
\nabla u(x^*)=\lambda p,\qquad p\cdot x^*=w.
\]
因而 $u_1(x^*)/u_2(x^*)=p_1/p_2$。
\end{theorem}
\begin{proof}
构造拉格朗日函数 $\mathcal L=u(x)+\lambda(w-p\cdot x)$。内点意味着非负约束不绑定，故一阶条件为
$u_i-\lambda p_i=0$。严格拟凹性保证满足条件的束为唯一全局最优解。
\end{proof}

\subsection{包络定理、Roy 恒等式与 Shephard 引理}
间接效用 $v(p,w)=u(x(p,w))$。在可微且唯一最优的情形，包络定理给出
\[
\frac{\partial v}{\partial w}=\lambda,\qquad
\frac{\partial v}{\partial p_i}=-\lambda x_i(p,w).
\]
消去 $\lambda$ 得 Roy 恒等式：
\[
x_i(p,w)=-\frac{\partial v(p,w)/\partial p_i}{\partial v(p,w)/\partial w}.
\]
对支出函数 $e(p,\bar u)=\min\{p\cdot x:u(x)\ge\bar u\}$，包络定理给出 Shephard 引理
\[
\frac{\partial e(p,\bar u)}{\partial p_i}=h_i(p,\bar u),
\]
其中 $h$ 是补偿需求。Roy 恒等式从最大化问题恢复马歇尔需求，Shephard 引理从最小化问题恢复希克斯需求；两者通过对偶关系相互校验。

\begin{example}[Cobb--Douglas 的双重校验]
令 $u(x_1,x_2)=x_1^\alpha x_2^{1-\alpha}$，$0<\alpha<1$。马歇尔需求为
$x_1=\alpha w/p_1$、$x_2=(1-\alpha)w/p_2$，间接效用为
\[
v(p,w)=\left(\frac{\alpha w}{p_1}\right)^\alpha
\left(\frac{(1-\alpha)w}{p_2}\right)^{1-\alpha}.
\]
对 $p_1$ 求偏导并代入 Roy 恒等式，恰好恢复 $x_1=\alpha w/p_1$。对偶问题的支出函数为
\[
e(p,\bar u)=\frac{\bar u}{\alpha^\alpha(1-\alpha)^{1-\alpha}}p_1^\alpha p_2^{1-\alpha},
\]
对 $p_1$ 求导则恢复 Hicks 需求。
\end{example}

\subsection{反例库：存在性条件为何不可省略}
\begin{itemize}
\item 若可行集不闭，连续效用的上确界可能无法达到，例如 $X=(0,1)$、$u(x)=x$。
\item 若效用不连续，最优选择可能随价格发生跳跃，比较静态不能直接使用导数。
\item 若偏好不拟凹，局部一阶条件可能只找到局部极值，不能保证全局最优。
\item 若价格向量含零分量，预算集可能无界，紧性证明必须另行处理。
\end{itemize}

\begin{remark}
书稿中的“存在性”与“计算性”必须区分：Weierstrass 定理保证最优解存在，但不告诉我们怎样算；一阶条件提供计算方程，但需要凹性或拟凹性把局部条件提升为全局结论。
\end{remark}

\subsection{综合证明与例题：存在性、识别与比较静态}

\subsection{一个完整的参数化例子}
考虑 $B(p,w)=\{x\in\mathbb R_+^2:p_1x_1+p_2x_2\le w\}$ 与
$u(x)=\min\{x_1,2x_2\}$。由于效用是凹函数，最优点位于 $x_1=2x_2$ 与预算线的交点，故
\[
x_1^*=\frac{2w}{p_1+2p_2},\qquad x_2^*=\frac{w}{p_1+2p_2},\qquad
v(p,w)=\frac{w}{p_1+2p_2}.
\]
该例展示：偏好结构先决定等效用线的形状，约束再决定可实现的切点。

\begin{figure}[H]\centering
\begin{tikzpicture}[x=0.9cm,y=0.9cm,font=\scriptsize]
\fill[gray!12] (0,0)--(0,3.8)--(4.2,0)--cycle;
\draw[->] (0,0)--(4.7,0) node[right] {$x_1$};\draw[->] (0,0)--(0,4.3) node[above] {$x_2$};
\draw[thick] (0,3.8)--(4.2,0) node[pos=.25,above left] {$p\cdot x=w$};
\draw[dashed] (0,2.1)--(4.2,0) node[pos=.52,below right] {$u=1.05$};
\draw[dashed] (0,1.15)--(2.3,0) node[pos=.45,below right] {$u=0.58$};
\fill (2.53,1.26) circle (2pt) node[above right] {$x^*$};
\end{tikzpicture}
\caption{完美互补偏好的最优点}
\end{figure}
\figexplan{两条虚线是 L 形偏好的等效用线的斜段示意，预算边界与最高可达等效用线在比例固定的点相交。最优解不是切线条件，而是固定比例条件与预算约束共同决定。}

\subsection{比较静态的符号判断}
由 $v=w/(p_1+2p_2)$，有 $\partial v/\partial w>0$、$\partial v/\partial p_i<0$。当 $p_1$ 上升时，$x_1^*$ 和 $x_2^*$ 均下降，原因是该例中两种商品严格互补。与之相比，替代品模型可能出现交叉价格效应为正；因此不能仅凭预算收缩判断每个分量的方向。

\subsection{存在性与唯一性的逻辑分离}
Weierstrass 定理只保证最大值被达到，不保证唯一性。严格拟凹性保证最优集合至多一个点；一般拟凹性允许一段最优点；非拟凹性甚至可能有多个分离的局部峰。研究写作中必须分别报告“存在”“唯一”和“可微”，三者是不同结论。

\begin{table}[H]\centering\small
\begin{tabular}{p{3cm}p{4cm}p{4cm}}
\toprule
条件 & 可推出什么 & 不能推出什么\\
\midrule
连续 $u$，紧 $B$ & 最优解存在 & 不保证唯一、可微\\
严格拟凹 & 最优解唯一（若存在） & 不保证闭式解\\
二阶可微 & 可用 Hessian 检验局部性质 & 不自动保证全局最优\\
严格凸约束集 & 几何结构更规整 & 不替代目标函数的凹性\\
\bottomrule
\end{tabular}
\caption{最优性结论的条件分层}
\end{table}

\subsection{综合练习}
\begin{enumerate}
\item 证明预算集在 $p\gg0$、$w<\infty$ 时紧。
\item 对 $u=x_1^\alpha x_2^{1-\alpha}$ 推导 Marshall 与 Hicks 需求，并验证 Slutsky 对称性。
\item 构造一个连续但非拟凹的效用函数，说明一阶条件可能产生多个候选点。
\item 解释为什么相同的选择预测并不意味着相同的福利判断。
\end{enumerate}

\subsection{现有模型的政策例题：参数、可行集与福利边界}

\subsection{税收与补贴的统一表示}
令一次性税收为 $T$、从价税率为 $\tau$，预算约束写为
$(1+\tau)p\cdot x\le w-T$。该写法把税收同时分解为相对价格效应与财富效应：$\tau$ 改变预算线斜率，$T$ 平行移动预算线。补贴则可表示为负税率或增加有效财富，取决于政策是否按数量给付。

\begin{figure}[H]\centering
\begin{tikzpicture}[x=0.9cm,y=0.8cm,font=\scriptsize]
\draw[->] (0,0)--(6,0) node[right] {$x_1$};\draw[->] (0,0)--(0,5) node[above] {$x_2$};
\draw[thick] (0,4.2)--(5.2,0) node[pos=.22,above left] {税前};
\draw[dashed] (0,3.1)--(5.2,0) node[pos=.2,above left] {对 $x_2$ 征税};
\draw[dotted] (0,4.2)--(4.1,0) node[pos=.65,below] {一次性税收};
\node[align=center,fill=white] at (2.0,3.4) {相对价格效应};\node[align=center,fill=white] at (4.6,1.6) {财富效应};
\end{tikzpicture}
\caption{税收的斜率效应与平移效应}
\end{figure}
\figexplan{实线为税前预算线，虚线改变斜率，点线保持斜率但向内平移。图中两个政策效应可以分别识别，也可在一般模型中同时存在。}

\subsection{补偿原则与效率判断}
若政策使某消费者效用下降，可用补偿变差衡量恢复到原效用所需的收入。福利分析不能只比较选择数量，还要明确基准价格、效用水平及补偿规则。帕累托改进要求没有人变差且至少一人变好；成本收益分析则允许有人受损，但要求总补偿可行。

\subsection{模型边界}
当预算集包含公共品、配额或非线性价格时，简单的线性预算式不再充分。应把制度规则直接写入可行集，而非事后把偏离归因于“偏好变化”。这也是本章从形式化基础进入政策应用的关键边界。

\subsection{现有方法的证明检查：从模型到可证伪命题}
\subsection{参数、机制与预测}
任何模型都应把参数分为偏好参数、技术参数和制度参数。参数不是“越多越好”：每增加一个自由参数，就必须说明其识别来源以及对预测的边际贡献。一个可证伪命题至少包含对象、参数变化、结果变量和方向判断。

\subsection{局部与全局}
局部比较静态围绕基准点展开，依赖导数和非奇异雅可比矩阵；全局结论则要求凹性、单调性或格理论等结构条件。二者不能混用。若论文只证明了局部结论，应明确写出“在给定邻域内”。

\subsection{研究设计检查表}
\begin{enumerate}
\item 可行集是否非空、闭且有界？\item 目标函数是否连续？\item 最优解是否唯一？\item 参数是否可由数据识别？\item 反事实是否超出支持集？\item 结论是否依赖某个不可检验的偏好限制？
\end{enumerate}
\subsection{讨论题}
请分别以价格补贴、配额和信息披露为例，写出可行集变化、最优集合变化以及福利判定所需的额外条件。每个案例均应区分正向预测与规范评价。

\subsection{模型价值的证明性讨论}
模型的有用性不仅取决于是否能够解释一个例子，还取决于它能否在新环境中给出稳定预测。准确性要求预测与观测相符；广泛性要求同一机制可迁移；简洁性要求不以无约束参数堆积换取拟合；识别性要求参数可由数据区分；可证伪性要求存在可能否定模型的事实。六者可能发生张力，研究者应在文中明确取舍。

\subsection{从反例学习}
反例不是理论失败的附录，而是界定适用范围的工具。若一个命题在非闭预算集上失败，应修正存在性条件；若在非凸偏好上失败，应修正唯一性结论；若在不可比关系上失败，应放弃单一序数效用的强断言。每个反例都对应一条必须补充的假设。

\subsection{综合案例：平台规则改变}
设平台规则使可行集从 $B_0$ 变为 $B_1$。正向问题是比较 $C(B_0)$ 与 $C(B_1)$；规范问题则要进一步比较效用、收入分配与外部性。不能把选择变化直接等同于福利变化，因为可行集扩张可能同时伴随风险上升或信息成本增加。

\subsection{本节练习}
\begin{enumerate}
\item 给出一个可行集扩张但最大值下降的例子，并说明需要什么效用定义。
\item 说明为什么选择对应比选择点更适合表示多重最优。
\item 对一个政策反事实分别写出正向命题、机制命题和规范命题。
\end{enumerate}

\subsection{现有定义的证明边界}
\subsection{三种证明模板}
存在性证明通常依次验证非空性、紧性与连续性；唯一性证明通常依次验证严格凸性或严格拟凹性；比较静态证明通常依次验证目标函数、约束和参数映射的正则性。把证明拆成模板有助于避免把不同层次的结论混在一起。

\subsection{边界条件的经济解释}
当价格趋近于零、收入趋近于零或配额趋近于上限时，内点条件可能失效。研究中应明确说明参数空间是否排除这些边界，以及结论能否连续延拓到边界。边界不是技术噪声，而是政策阈值和制度约束最可能出现的区域。

\subsection{小结}
本章的形式化主线可以压缩为：先规定可行性，再规定评价，再定义选择，最后检验预测。任何跳过其中一步的论证，都可能把制度约束、偏好判断与经验事实混为一谈。

\subsection{核心结论的逐步证明与反例核对}

\begin{proposition}[预算集的逐坐标界]
若 $p\gg0$ 且 $x\in B(p,m)$，则对任意 $i$ 有 $0\le x_i\le m/p_i$。
\end{proposition}
\begin{proof}
由 $x\in\R_+^L$ 得 $x_i\ge0$。另一方面，$p_jx_j\ge0$ 对每个 $j$，故
$p_ix_i\le\sum_jp_jx_j=p\cdot x\le m$。由于 $p_i>0$，两边除以 $p_i$ 得 $x_i\le m/p_i$。这一逐坐标界同时给出预算集有界性，并说明价格严格为正是有界性不可缺少的条件。
\end{proof}

\begin{proposition}[最大值定理的应用]
若 $B$ 非空紧，$u:B\to\R$ 连续，则 $\argmax_{x\in B}u(x)$ 非空且紧。
\end{proposition}
\begin{proof}
由 Weierstrass 定理，存在 $x^*\in B$ 使 $u(x^*)\ge u(x)$ 对所有 $x\in B$，故最优集合非空。令 $M=u(x^*)$，则
$\argmax_Bu=B\cap u^{-1}(\{M\})$。单点集 $\{M\}$ 在实数中闭，连续性使其原像闭；闭集与紧集的交仍紧，故结论成立。
\end{proof}

\begin{example}[去掉闭性后的失败]
取 $B=(0,1)$、$u(x)=x$。函数连续且 $B$ 有界，但对每个 $x<1$ 都有 $(x+1)/2\in B$ 且 $u((x+1)/2)>u(x)$，因此不存在最大元。若改为 $B=[0,1]$，同一函数在 $x=1$ 处达到最大值。失败来自集合的边界被删除，而非来自效用函数不连续。
\end{example}

\begin{example}[去掉有界性后的失败]
取 $B=[0,\infty)$、$u(x)=x$。集合闭但无界，且对任意候选点 $x$ 都存在 $x+1$ 使效用更高。该例说明“闭”只保证极限点不丢失，不能阻止改进序列逃向无穷远。
\end{example}

\begin{proposition}[向下闭合与财富耗尽]
若 $x\in B(p,m)$、$0\le\lambda\le1$，则 $\lambda x\in B(p,m)$。若 $u$ 局部非饱和且 $m>0$，任何最优点都满足 $p\cdot x=m$。
\end{proposition}
\begin{proof}
第一结论已由 $p\cdot(\lambda x)=\lambda p\cdot x\le\lambda m\le m$ 得出。若某最优点满足 $p\cdot x<m$，则存在 $\varepsilon>0$ 使一个足够小的正方向仍留在预算集内。局部非饱和意味着沿某个正方向可得到严格更高效用，与最优性矛盾。因此预算约束绑定。
\end{proof}

\subsection{比较静态证明的逐步展开}
设最优束由方程 $F(x,\theta)=0$ 定义，且 $F_x$ 在基准点可逆。对恒等式求导：
\[
F_x(x(\theta),\theta)\frac{\partial x}{\partial\theta_k}+F_{\theta_k}(x(\theta),\theta)=0.
\]
于是
\[
\frac{\partial x}{\partial\theta_k}=-F_x^{-1}F_{\theta_k}.
\]
该公式的每一个符号都有明确含义：$F_x^{-1}$ 描述最优条件对状态变量的敏感度，$F_{\theta_k}$ 描述参数对最优条件的直接冲击。只有两者都定义良好，比较静态方向才有数学依据。

\begin{example}[两商品 Cobb--Douglas 的逐步推导]
令 $u=x_1^\alpha x_2^{1-\alpha}$，预算约束为 $p_1x_1+p_2x_2=m$。拉格朗日函数为
\[
\mathcal L=x_1^\alpha x_2^{1-\alpha}+\lambda(m-p_1x_1-p_2x_2).
\]
一阶条件为
\[
\alpha x_1^{\alpha-1}x_2^{1-\alpha}=\lambda p_1,
\quad (1-\alpha)x_1^\alpha x_2^{-\alpha}=\lambda p_2.
\]
两式相除，得到
\[
\frac{\alpha}{1-\alpha}\frac{x_2}{x_1}=\frac{p_1}{p_2},
\quad\text{即}\quad x_2=\frac{1-\alpha}{\alpha}\frac{p_1}{p_2}x_1.
\]
代回预算约束：
\[
p_1x_1+p_2\frac{1-\alpha}{\alpha}\frac{p_1}{p_2}x_1
=\frac{p_1}{\alpha}x_1=m,
\]
所以 $x_1=\alpha m/p_1$，再代回得 $x_2=(1-\alpha)m/p_2$。这一步同时验证了一阶条件、预算耗尽和需求的齐次性。
\end{example}

\subsection{对偶关系的证明链}
定义 $v(p,m)=\max_{p\cdot x\le m}u(x)$。若最优解唯一且可微，包络定理给出
$v_m=\lambda$、$v_{p_i}=-\lambda x_i$。因此
\[
-\frac{v_{p_i}}{v_m}=x_i.
\]
这不是额外假设，而是把同一个最优问题分别对财富和价格求导后消去乘子。支出问题
$e(p,\bar u)=\min\{p\cdot x:u(x)\ge\bar u\}$ 的包络式为 $e_{p_i}=h_i$，其中 $h_i$ 是补偿需求。两条恒等式必须在同一组偏好与正则条件下使用。

\subsection{证明细节：从集合性质到选择结论}

\begin{proof}[Heine--Borel 条件在预算集上的逐步验证]
对任意序列 $x^n\in B(p,m)$，由逐坐标界 $0\le x_i^n\le m/p_i$，序列落在有界盒 $\prod_i[0,m/p_i]$ 内。Bolzano--Weierstrass 定理给出收敛子列 $x^{n_k}\to x^*$。非负性由坐标极限保持，且连续性给出
$p\cdot x^*=\lim_k p\cdot x^{n_k}\le m$，故 $x^*\in B(p,m)$。这直接验证了紧性的序列刻画，不必只引用闭且有界的结论。
\end{proof}

\begin{proof}[最优集合为何是闭集]
令 $M=\max_{x\in B}u(x)$，取任意序列 $x^n\in\argmax_Bu$ 且 $x^n\to x$。由于每个 $u(x^n)=M$，连续性给出 $u(x)=\lim_nu(x^n)=M$。若 $B$ 闭，则 $x\in B$，所以 $x\in\argmax_Bu$。这说明最优集合不仅非空，而且对极限过程稳定。
\end{proof}

\begin{example}[最优集合不必是单点]
令 $B=\{x\in\R_+^2:p\cdot x\le m\}$，$u(x)=p\cdot x$。所有预算边界点都达到 $m$，因此 $C(p,m)=\{x:p\cdot x=m,x\ge0\}$ 是一条线段。连续性和紧性保证这条线段存在，但只有严格拟凹性才能排除多重最优。
\end{example}

\subsection{一阶条件与约束资格的核对}
在内点情形，拉格朗日一阶条件是必要条件；在角点情形，必须改用 Kuhn--Tucker 条件。对
$\max u(x)$ s.t. $g(x)\le0$、$x\ge0$，拉格朗日函数可写为
\[
\mathcal L=u(x)-\lambda g(x)+\mu\cdot x,
\qquad \lambda\ge0,\ \mu_i\ge0.
\]
其条件包括原始可行性、对偶可行性、互补松弛和驻点条件。内点时 $\mu=0$，才退化为边际效用比例等于价格比例。

\begin{example}[角点与内点的区别]
若 $u(x)=x_1+\frac12x_2$、预算为 $x_1+x_2\le m$，商品一的单位效用为 $1$，商品二为 $1/2$，最优是 $(m,0)$。若机械套用内点切线条件，会错误地要求 $1=1/2$。正确判断是比较单位价格下的边际效用，并检查非负约束的乘子。
\end{example}

\subsection{对偶恒等式的再推导}
由 $v(p,m)=u(x(p,m))$，全微分为
$dv=\nabla u(x)\cdot dx$。在内点最优点有 $\nabla u=\lambda p$，而预算恒等式给出
$p\cdot dx+x\cdot dp=dm$。故
\[
dv=\lambda(dm-x\cdot dp),
\]
从而 $v_m=\lambda$、$v_{p_i}=-\lambda x_i$。这给出 Roy 恒等式，也说明价格导数为负并非独立假设，而是最优性与预算恒等式共同推出的结果。

\begin{remark}
本章所有比较静态结论均应回到三项检查：最优解是否存在、是否唯一、是否位于内点。只有三项条件分别核对后，导数符号才具有清晰的经济含义。
\end{remark}

\subsection{证明细节：选择对应、最大值与比较静态}

\begin{proposition}[选择对应的集合值性质]
设 $C(\theta)=\argmax_{x\in B(\theta)}u(x)$。若 $B(\theta)$ 非空紧、$u$ 连续，则 $C(\theta)$ 非空闭，并且 $C(\theta)\subseteq B(\theta)$。
\end{proposition}
\begin{proof}
包含关系是 argmax 定义的直接结果。非空性由最大值定理得到。若 $x^n\in C(\theta)$ 且 $x^n\to x$，因为 $B(\theta)$ 闭，$x\in B(\theta)$；又因 $u(x^n)=M$ 且 $u$ 连续，$u(x)=M$，故 $x\in C(\theta)$。因此选择对应是闭集，而不是任意的点集。
\end{proof}

\begin{example}[线性效用的整段最优]
令 $u(x)=x_1+2x_2$，预算线为 $p_1x_1+p_2x_2=m$。若 $p_1/p_2=1/2$，预算线与效用等值线平行，所有预算边界点都是最优点；若 $p_1/p_2<1/2$，单位支出带来的商品一效用更高，最优为商品一角点；若 $p_1/p_2>1/2$，最优为商品二角点。三种情形分别对应整段、左端点和右端点，不能仅用一个“切点”公式概括。
\end{example}

\begin{proposition}[严格拟凹性与唯一性]
若 $B$ 为凸集、$u$ 严格拟凹，则 $\argmax_Bu$ 至多包含一个点。
\end{proposition}
\begin{proof}
反设 $x\ne y$ 都是最大点，记最大值为 $M$。凸性给出 $z=\lambda x+(1-\lambda)y\in B$，其中 $0<\lambda<1$。严格拟凹性给出
$u(z)>\min\{u(x),u(y)\}=M$，与 $M$ 是最大值矛盾。因此最优点至多一个。
\end{proof}

\subsection{参数变化的逐步判别}
预算集 $B(p,m)$ 随收入 $m$ 增加而扩张，随价格上升而收缩，但最优集合的变化还取决于效用函数。若 $u$ 严格单调，最优点位于边界；若 $u$ 存在饱和点，收入增加可能不改变选择。比较静态必须同时写出约束变化和目标函数结构，不能从集合扩张直接推出每个消费量上升。

\begin{example}[收入扩张与非正规品]
若效用为 $u(x_1,x_2)=x_1-(x_2-2)^2$，收入增加首先使 $x_2$ 向饱和点 $2$ 调整，超过该点后继续增加收入并不必然增加 $x_2$。这说明“收入上升、需求上升”是正规品假设下的结论，而非预算集扩张的逻辑必然。
\end{example}

\begin{proposition}[包络公式的直接证明]
令 $v(\theta)=u(x(\theta),\theta)$，且 $x(\theta)$ 为内点唯一最优解、满足一阶条件 $u_x=0$。则
$v_{\theta}=u_{\theta}(x(\theta),\theta)$。
\end{proposition}
\begin{proof}
链式法则给出 $v_{\theta}=u_x(x(\theta),\theta)x_{\theta}+u_{\theta}(x(\theta),\theta)$。最优一阶条件令第一项为零，故只剩直接参数效应。约束问题中把约束乘子写入拉格朗日函数后，同样的消去机制给出价格和收入的包络导数。
\end{proof}

\subsection{模型边界的三个反例}
\begin{enumerate}
\item 若 $p_i=0$，预算集沿商品 $i$ 方向可能无界，紧性论证失效。
\item 若 $u$ 不连续，闭预算集也不能保证最大值达到；跳跃效用可能只有上确界。
\item 若可行集不凸，一阶条件即使存在，也可能只识别某一局部峰值。
\end{enumerate}
上述反例不是新增模型，而是对本章既有存在性、唯一性和一阶条件结论的边界核查。写作时应把每个定理的假设与结论逐项对应，避免把一个结论的假设带入另一个结论。

\subsection{同一选择问题的三种边界情形}
设预算集为 $B(p,m)=\{x\in\mathbb R_+^2:p_1x_1+p_2x_2\le m\}$。对线性效用
$u(x)=a_1x_1+a_2x_2$，比较单位支出的效用收益：
\[
\rho_1=\frac{a_1}{p_1},\qquad \rho_2=\frac{a_2}{p_2}.
\]
若 $\rho_1>\rho_2$，最优集合退化为商品一角点；若 $\rho_1<\rho_2$，最优集合退化为商品二角点。只有 $\rho_1=\rho_2$ 时，预算边界上的每一点都给出相同效用，因而
\[
C(p,m)=\{x\ge0:p\cdot x=m\}.
\]
这三个区域说明，角点、整段和唯一内点并不是互相矛盾的结论，而是同一最大化问题在不同参数区域的完整划分。

\begin{proposition}[财富耗尽的条件]
若 $u$ 对每个商品都严格单调、$p\gg0$ 且 $m>0$，则任意 $x^*\in C(p,m)$ 都满足 $p\cdot x^*=m$。
\end{proposition}
\begin{proof}
反设 $p\cdot x^*<m$。由于 $p_i>0$，可取充分小的 $\delta>0$，使得
$p\cdot(x^*+\delta e_i)\le m$。严格单调性给出
$u(x^*+\delta e_i)>u(x^*)$，这与 $x^*$ 已达到可行集上的最大值矛盾。因此最优点必须位于预算边界。若效用只满足弱单调，财富耗尽则不再自动成立。
\end{proof}

\subsection{从候选点到全局最优：一项完整核对}
设 $B=\{x\ge0:p\cdot x\le m\}$，且 $u$ 连续、严格递增并严格拟凹。首先由 $p\gg0$ 验证 $B$ 有界，由闭性验证 $B$ 闭，故 $B$ 紧；最大值定理给出最优解存在。其次，严格递增性排除预算内部点，故 $p\cdot x^*=m$。最后，严格拟凹性排除两个不同的最大点，因而最优解唯一。

这一链条中每一步承担不同任务：紧性只负责“能取到最大值”，严格单调只负责“财富耗尽”，严格拟凹只负责“不能有两个最优点”。若删除其中任一条件，结论都必须相应减弱。比如连续但不严格拟凹的线性效用可以产生整段最优；严格拟凹但预算集不紧时，可能只有上确界而没有最大值。

\begin{example}[两项假设的逐项删去]
取 $u(x)=x_1+x_2$，它连续但不是严格拟凹，预算线上所有点都最优；若预算集改为开集 $\{x\ge0:p\cdot x<m\}$，即使 $u$ 连续且严格递增，也只能不断逼近边界，最大值可能不存在。这两个例子分别对应唯一性与存在性的失败。
\end{example}

\subsection{候选点、边界点与全局比较的统一算法}
对于本章反复出现的约束最大化问题，可以把证明拆成四个互不替代的步骤。第一步写出可行集并逐项检查非空、闭性和有界性；第二步利用连续性得到最大值存在；第三步根据单调性判断最优解是否耗尽约束；第四步才使用一阶条件、Kuhn--Tucker 条件或无差异曲线进行定位。最后必须将候选点与全部边界点比较，才能把局部条件升级为全局结论。

例如，对 $u(x)=\alpha\log x_1+(1-\alpha)\log x_2$，定义域要求 $x_1,x_2>0$。在正象限内，拉格朗日一阶条件给出
\[
\frac{\alpha}{x_1}=\lambda p_1,\qquad
\frac{1-\alpha}{x_2}=\lambda p_2.
\]
两式相除得到 $x_1/x_2=\alpha p_2/[(1-\alpha)p_1]$，再代入预算等式得到候选束。由于对数效用严格凹，任意满足一阶条件的候选束都是唯一全局最优；这里“严格凹”承担的是充分性和唯一性，而不是存在性本身。

若把效用换成线性函数，上述一阶条件可能没有内点解，必须比较两个角点和预算线整段。若把预算集换成不闭的开集，连续性也不再足以保证最大值达到。由此可见，证明顺序不是形式上的排版要求，而是防止把候选解误写成最优解的逻辑防线。

\subsection{一阶条件的充分性为何需要曲率条件}
设 $x^*$ 是内点候选解，满足 $\nabla u(x^*)=\lambda p$。这一等式只说明沿预算超平面的局部可行方向，效用的一阶变化为零，并不能单独排除另一个更高的可行点。若 $u$ 凹，则任意 $y\in B(p,m)$ 满足
\[
 u(y)\le u(x^*)+\nabla u(x^*)\cdot(y-x^*)
 =u(x^*)+\lambda p\cdot(y-x^*)\le u(x^*),
\]
因此候选点是全局最优。若 $u$ 严格凹，最后一个不等式在 $y\ne x^*$ 时严格，故最优唯一。这个两行不等式正是“局部一阶条件加全局曲率”升级为全局结论的关键。

若只有拟凹性而没有可微性，则不能直接写梯度条件，但仍可用上水平集的凸性和支持超平面描述最优点；若有可微性而没有凹性，则一阶条件只给出驻点，必须回到边界比较或二阶条件。因而“写出一阶条件”与“证明最优性”在逻辑上始终是两件事。

\subsection{本章证明链的回顾与应用边界}
从方案空间到选择对应，完整论证依次经过集合的定义、可行性的限制、最大值的存在、最优集合的识别以及参数变化的比较。预算集的紧性保证“最大值能够取到”，连续性把集合性质传递给目标函数，单调性决定边界位置，严格拟凹性决定唯一性，而微分条件只是在这些前提成立后提供计算工具。任何一步被省略，结论都只能保持在较弱层次。

因此，在阅读或书写消费者问题时，应同时记录“结论—假设—证明工具”三列。例如，存在性对应紧性与连续性；财富耗尽对应严格单调性与正价格；唯一性对应严格拟凹性；需求导数对应内点、可微性和隐函数定理。这个核对表既能解释定理为何成立，也能在参数改变时迅速指出哪一条结论需要重新证明。

\begin{example}[同一模型中的假设替换]
若把严格单调替换为弱单调，预算耗尽结论必须改写为“至少存在一个边界最优解”，不能再断言每个最优解都在边界；若把严格拟凹替换为普通拟凹，存在性仍可保留，但唯一性消失；若把闭预算集替换为开预算集，连续性仍可保留，但最大值定理不再适用。逐项替换说明，假设不是装饰性的前置条件，而是分别支撑不同结论的逻辑支柱。
\end{example}

\subsection{参数、集合与选择的三层区分}
设参数从 $\theta$ 变为 $\theta'$，首先要比较的是可行集 $B(\theta)$ 与 $B(\theta')$；其次才比较目标函数在两个集合上的最大值；最后比较选择对应 $C(\theta)$ 与 $C(\theta')$。即使 $B(\theta)\subseteq B(\theta')$，也只能说明可行机会扩大，不能单独推出每一坐标需求增加，更不能推出福利必然增加，除非再加入单调性、补偿或规范评价的额外条件。

在唯一最优的情形，选择对应可写成函数 $x(\theta)$，可以讨论连续性和导数；在多重最优的情形，必须保留集合值记号 $C(\theta)$，并使用上半连续、下半连续或选择函数来描述变化。把集合值对应强行写成单值函数，会掩盖线性效用、完全替代和约束退化时的真实结构。该区分也是本章从集合语言出发、最后回到模型价值评价的闭环。


% ==================== 参考文献 ====================
\clearpage
\section*{本章参考文献}
\begin{enumerate}
  \item Arrow, K. J. (1951). \emph{Social Choice and Individual Values}. Yale University Press.
  \item Debreu, G. (1959). \emph{Theory of Value}. Wiley.
  \item Afriat, S. N. (1967). The construction of separable utility functions from expenditure data. \emph{International Economic Review}, 8(1), 67–77.
  \item Sen, A. (1970). \emph{Collective Choice and Social Welfare}. Holden-Day.
  \item Varian, H. R. (1992). \emph{Microeconomic Analysis} (3rd ed.). Norton.
  \item Hausman, D. M. (1992). \emph{The Inexact and Separate Science of Economics}. Cambridge University Press.
  \item Mas-Colell, A., Whinston, M. D., \& Green, J. R. (1995). \emph{Microeconomic Theory}. Oxford University Press.
\end{enumerate}

\end{document}
